\section{The notes of 1983: centralizers, descent, and tangential parameters}
\label{sec:F}

We write $G_K=\operatorname{Gal}(\overline K/K)$,
$H_K=\ker(\chi_K:G_K\to\widehat{\mathbb Z}^{\times})$,
$Q_K=\chi_K(G_K)$, $T=\widehat{\mathbb Z}(1)$,
$P_U=\pi_1(U_{\overline K})$, and $E_U=\pi_1(U)$.
The notation $\widehat{A}$ for an abelian group means
$\varprojlim_n A/nA$, or $\varprojlim_n A/A^n$ in multiplicative
notation. Its terms need not be finite. All section classes are
classes under geometric inner conjugation. An actual section, an
actual lifted cusp, and an actual subgroup are specified whenever a
centralizer is computed. The comparisons in this section concern cote
149, scans 44--63, 85, 100--103, and 126--127. We retain the two written
subscripts $L$ on scan 127 and separately state their mathematical
correction.

\subsection{F(a)--(c): what the centralizer calculations prove}

\begin{proposition}[F.1 --- Known, with checked reduction: the cyclotomic centralizer]
\label{F-centralizer}
Let $K/\mathbb Q$ be finitely generated and let $U/K$ be a smooth
geometrically connected hyperbolic curve. For every continuous section
$s:G_K\to E_U$,
\[
 C_{P_U}(s(H_K))=1
 \quad\hbox{or}\quad
 C_{P_U}(s(H_K))=I_c\simeq T,
\]
where the second case holds precisely when $s(G_K)$ is contained in
the decomposition group of a unique lifted rational cusp $c$.
The centralizer of a point section is $1$. Restricting the section to
an open subgroup of $H_K$ preserves this centralizer. This gives the
two ``pr\'ecisions'' on cote 149, scan 56, and the formula on scan 55.
More generally, for a closed subgroup $B\subset E_U^{\circ}$
whose projection to $H_K$ is open, $C_{P_U}(B)$ is trivial or
one entire cusp inertia. Consequently the same proof recognizes
every nontrivial procyclic subgroup essentially cyclotomically
normalized over an open subgroup of $G_K$; no splitting of that
normalizer is required.
\end{proposition}
\begin{proof}
We checked CYCLOTOMIC\_KERNEL A1--A7 and AE1. We record the steps
which carry the arithmetic assertion to the full profinite group;
in particular, we do not discard its common kernel of pro-prime
quotients. Ribet's Appendix, Theorem 1, gives finite torsion in
$A(k(\mu_\infty))$ for an abelian variety over a number field
$k$~\cite{F-Ribet}. For a finitely generated $K/k$, spread $A$ to an
abelian scheme over a smooth integral $k$-variety $S$. Each rational
torsion point over $Kk(\mu_\infty)$ extends over $S_{k(\mu_\infty)}$:
the appropriate torsion scheme is finite \'etale, and its generic
degree-one component extends by normality. Evaluating these sections
at a closed point $x$ gives an injection into
$A_x(k(x)(\mu_\infty))_{\rm tors}$. Equality at one fibre is equality
everywhere since the equalizer of two sections of a finite \'etale
scheme is open and closed. Thus every compactified covering Jacobian
has zero $H_K$-invariants in its full Tate module.

For $z\in C_{P_U}(s(H_K))$, and every open characteristic
$N\subset P_U$, put $J=N\overline{\langle z\rangle}$. An open
subgroup of $G_K$ stabilizes $J$, and $Js(G_{K'})$ defines the
corresponding finite cover. The image of $z$ in $J^{\rm ab}$ maps
to zero in the Tate module of its compactification. The localization
sequence is
\[
 T\longrightarrow\bigoplus_{d\in D_J} T e_d
 \longrightarrow J^{\rm ab}\longrightarrow T\operatorname{Jac}(X_J)
 \longrightarrow0,
\]
with the diagonal as the first map. In the proper case the cyclic
quotient $J/N$ then forces $z=1$. In the affine case its image in
$J/N$ belongs to the subgroup generated by cusp inertia images.
As $zN$ generates the finite cyclic group $J/N$, those images generate
all of $J/N$. For a prime $\ell$, subgroups of a finite cyclic
$\ell$-group are linearly ordered. Hence one of these inertia images
contains the full $\ell$-Sylow subgroup of $J/N$, and therefore
contains the image of the $\ell$-primary component $z_\ell$ of $z$.
This inertia image is the stabilizer of a cusp in the normal cover
with subgroup $N$. Thus $z_\ell$ fixes a cusp at this level.
For any finite collection of levels a common normal refinement
gives compatible fixed cusps, so compactness gives a lifted cusp
fixed by $z_\ell$. This is the finite-cyclic argument of CK A3;
no assertion that an arbitrary sum of boundary classes belongs
to one boundary line is used. Distinct cusp inertias
intersect trivially, and the centralizer of a nontrivial element of
one inertia is that inertia (Stix, Proposition 11 and Corollary
12~\cite{F-Stix}). Since the primary components commute, they lie in
one inertia: fix any nonzero primary component, and use its
inertia as its geometric centralizer to contain every other
component. The element $z$ is the limit of its finite products
of primary components (in each finite cyclic quotient only
finitely many primes occur). The inertia is closed, so it
contains $z$. Thus compactness across covering levels and the
comparison across different primes are separate justified steps.

Here is the necessary passage from elements to the entire fixed
subgroup, checked in CK A5. A closed subgroup all of whose nonidentity
elements lie in inertia lies in a single inertia. Otherwise choose
two distinct orbits of its cusps: a finite group cannot be the union
of conjugates of one proper subgroup, by counting fixed cosets.
They remain distinct in a finite cover containing the subgroup.
Refine to at least three cusps; when there are exactly two, the
compactification has positive genus and an \'etale double cover of
it supplies four. More explicitly, if $C$ is the fixed subgroup,
choose the two cusps in distinct $C$-orbits and an open normal
$N$ at which these orbits are still distinct. Such an $N$ exists
by compactness of $C$. The cover with subgroup $CN$ contains $C$
and separates the two cusps. Every element of $C$ maps to the
identity in the fundamental group of its compactification,
since it is an inertia element in this cover. Hence the pullback
of the stated compact double cover still contains $C$; no
arbitrary refinement containing $C$ is being presumed.
In the boundary module, take nonzero images
$ae_c,be_{c'}$. Both products giving $ae_c+be_{c'}$ and
$2ae_c+be_{c'}$ would have to lie in individual boundary lines.
With at least four cusps neither can do so; with three the two
requirements imply $a=b$ and $2a=b$, hence $a=0$. This contradiction
works over $\widehat{\mathbb Z}$ without dividing by a zero divisor.
We obtain $C\subset I_c$. A nonidentity $z\in C\cap I_c$ belongs
also to every $s(h)I_cs(h)^{-1}$ for $h\in H_K$. Trivial intersection
of distinct cusp inertias therefore makes $s(H_K)$ normalize $I_c$,
and the kernel of the cyclotomic character acts trivially on it;
therefore $I_c\subset C$. Normality of $H_K$ in $G_K$ makes the whole
section normalize $I_c$, giving the asserted rational cusp.

The separation from point sections uses the Jacobian Kummer map
after a cover whose compactification has positive genus. Here is
the precise divisible-point argument. First let $L=k(\mu_\infty)$
for a number field $k$, and descend an abelian variety $A/L$ and
a divisible point $P$ to a number field $k'\subset L$ with
$L=k'(\mu_\infty)$. The finite nonempty sets of factorial-division
points of $P$ in $A(L)$ admit a compatible inverse limit. Its
Kummer cocycle $c:G_{k'}\to T(A)$ vanishes on the normal subgroup
$H=G_L$. For $h\in H$, $g\in G_{k'}$, the cocycle identity evaluated
at $hg=g(g^{-1}hg)$ gives $h c(g)=c(g)$. The torsion calculation
above gives $T(A)^H=0$, so $c=0$. All division points are already
defined over $k'$. Reducing at one place of good reduction modulo
every power of its maximal ideal sends $P$ to a divisible element
of a finite group, hence to zero. These reductions separate points,
so $P=0$.

For finitely generated $K$, every finite extension of $K(\mu_\infty)$
is $K'(\mu_\infty)$ for one finitely generated $K'/\mathbb Q$.
Descend $A$ and a proposed nonzero divisible $P$ to $K'$, choose
a smooth geometrically integral model $S$ over its number field of
constants, and spread them to an abelian scheme and a section.
Shrink $S$ once so that this section avoids the identity section.
For every $n$, the division scheme $[n]^{-1}(P)$ is then finite
\'etale over this same $S$; no further shrinking depends on $n$.
Its rational point over $K'(\mu_\infty)$ extends over
$S_{k(\mu_\infty)}$ by normality. Evaluation at any closed point
of $S$ makes a nonzero abelian-variety point over a number-field
cyclotomic extension divisible by every $n$, contradicting the
preceding paragraph. This supplies exactly the hypothesis in
Stix's Theorem 14 over every finite extension of $K(\mu_\infty)$.
For a positive-genus curve, $[x-y]=0$ implies $x=y$, since a
function with divisor $x-y$ for distinct points has degree one.
The theorem therefore separates all point and cusp
packets. Finite extensions of the cyclotomic field have the form
$K'(\mu_\infty)$; the identical argument, followed by normality,
proves the open-subgroup assertion.

For the more general $B$, let $z\in C_{P_U}(B)$ and use the same
$J=N\overline{\langle z\rangle}$. Descend its finite geometric
cover to some finite extension $K'/K$, and let $E_J\subset E_U$
be its arithmetic open subgroup, with $E_J\cap P_U=J$.
The group $B\cap E_J$ is open in $B$ and its projection is an
open subgroup of $H_K$. Its action on $J^{\rm ab}$ factors through
that projection, because its geometric kernel lies in $J$ and
acts by inner automorphisms. The image of $z$ in the compactified
Jacobian is therefore fixed by an open subgroup of $H_K$, and
vanishes by the already proved arithmetic assertion. The cyclic
and boundary-module arguments apply unchanged and give a unique
inertia containing every centralizing element. The cyclotomic
character of $B$ is trivial, so that entire inertia centralizes $B$.
Finally, for a subgroup $N$ which projects openly to $G_K$ and
normalizes a nontrivial procyclic $A$ through $\chi_K$, set
$B=N\cap E_U^{\circ}$. Its projection is open in $H_K$ and it
centralizes $A$. Hence $A$ is contained in the recognized inertia.
This also supplies the exact comparison with Stix's essential
cyclotomic-normalizer criterion without inserting a nonexistent
normalizer section.
\end{proof}




\begin{proposition}[F.2 --- Proved: the necessary half of (F1)]
\label{F-geometric-necessary}
For an elementary hyperbolic polycurve $U/K$ over a number field and
any morphism $f:V\to U$ from a nonempty connected finite-type
$\mathbb Q$-scheme, every open subgroup $J\subset E_V^{\circ}$
satisfies
\[
 C_{P_U}(f_*(J))=C_{E_U}(f_*(J))=1.
\]
Here $E_V^{\circ}$ is the cyclotomic kernel, and $f_*$ denotes a
continuous representative formed using compatible geometric base
paths. Thus regularity of $V$ is unnecessary for this direction of
scan 48. Both readings of its word ``centralisateur'' give the same
condition.
\end{proposition}
\begin{proof}
For a rational point of $U$, induct along its tower of relative
hyperbolic curves. An element centralizing its cyclotomic section
maps to $1$ in the preceding stage, and then belongs to the
geometric fibre group. The fibre group injects and its point action
is the ordinary fibre point action: these are precisely the maps
in Hoshi, Proposition 2.4(i),(ii)~\cite{F-HoshiPoly}, and in CK C5.
Proposition~\ref{F-centralizer} kills the element at the last stage.
This checks the induction in CK C6, including its left exactness.

Choose open normal $N\subset E_V$ with $N\cap E_V^{\circ}\subset J$.
The connected finite \'etale cover $V_N$ has a closed point $v$ over
a number field $k$. Its cyclotomic point homomorphism lands in $J$,
and its composite under $f_*$ is the point section of $f(v)$. The
previous paragraph proves the geometric equality. An arithmetic
centralizer element projects to an element of $G_{\mathbb Q}$
centralizing $G_{k(\mu_\infty)}$. This centralizer is trivial:
replace $k$ by its Galois closure and apply Minamide--Tsujimura,
Theorem 2.12 and Proposition 1.2~\cite{F-MinamideTsujimura} to that
nontrivial closed normal subgroup. Its nontriviality follows, for
example, by choosing a completely split prime of the finite
coefficient field and an odd prime degree ramification not occurring
in its cyclotomic closure; CK B1 gives the complete valuation
argument. Thus the element lies in $P_U$ and is already trivial.
Every change of base paths conjugates the displayed groups and
preserves the asserted equality.
\end{proof}

\begin{proposition}[F.3 --- Proved: the point case and the cuspidal exclusion]
\label{F-point-equivalence}
For $V=\operatorname{Spec}K'$ and a hyperbolic curve $U/K$, condition
(F1)(b) holds exactly for noncuspidal section classes over $K'$, for
each finite extension $K'/K$. The point case of (F1) and (F2) are
therefore the same surjectivity assertion, with the same injective
point map. Every homomorphism through a cusp decomposition group
fails (F1)(b). Both point sections and cusp sections have infinite
index image in $E_U$.
\end{proposition}
\begin{proof}
Proposition~\ref{F-centralizer}, including its open-subgroup clause,
identifies (b) with absence from all cusp packets. Removing those
packets from the set asserted to be exhausted by (F2) leaves exactly
the image sought in the point case of (F1). The equality is of sets
and maps, not an appeal to either surjectivity claim. For a map into
$D_c=I_c\rtimes G_K$, its cyclotomic-kernel image acts trivially on
$I_c$, so the latter is a nonzero subgroup of its centralizer.
After choosing a parameter at $c$, a unit $u\in\mathcal O(V)^*$
gives such a map: its compatible Kummer torsors give a continuous
cocycle $a:E_V\to T$, and the homomorphism is
$h\mapsto(a(h),p(h))\in T\rtimes G_K\hookrightarrow E_U$.
A point or cusp section has trivial intersection with $P_U$;
the quotient of $E_U$ by its image has the infinite coset space
$P_U$. Openness therefore removes both kinds of point-source map.
\end{proof}

\begin{proposition}[F.4 --- Known, with reduction: open maps to a curve]
\label{F-open-Hom}
For a normal geometrically connected variety $V/K$, a hyperbolic
curve $U/K$, and a number field $K$, continuous open homomorphisms
$E_V\to E_U$ over $G_K$, modulo conjugation by $P_U$, are precisely
the homomorphisms of unique dominant $K$-morphisms $V\to U$.
Each such open homomorphism satisfies (F1)(b).
\end{proposition}
\begin{proof}
This is the ordinary finite-type specialization of Mochizuki,
Theorem 16.5 and its following remark, as checked in F149-ZCF1
and in the attached source~\cite{F-Mochizuki}. For an open map $f$,
the map $E_V^{\circ}\to f(E_V)\cap E_U^{\circ}$ is onto and open.
The image of every open subgroup therefore contains an open
subgroup of $P_U$. Its geometric centralizer is trivial by the
slimness of a hyperbolic curve group. Projection to $G_K$ and the
arithmetic centralizer calculation in
Proposition~\ref{F-geometric-necessary} give the same conclusion in
$E_U$. The groupoid comparison in Part A sends a morphism with its
base paths to precisely this outer homomorphism. Its degree on
inertia is the ramification index of the reconstructed morphism,
so the peripheral/Kummer data are the data of that same morphism.
\end{proof}

\begin{proposition}[F.5 --- Refuted: open image forces dominance in higher dimension]
\label{F-hyperplane}
The open-image dominance assertion for higher-dimensional elementary
targets is false. For $C:y^2=x^5-x$ over $\mathbb Q$ and $Y=C^3$,
let $H$ be the inverse image in $Y$ of the trilinear divisor
\[
 1+7x_1+3x_2+2x_3+13x_1x_2+11x_1x_3
       +5x_2x_3+19x_1x_2x_3=0
\]
in $(\mathbb P^1)^3$. The closed immersion $H\hookrightarrow Y$
is nondominant and induces $E_H\simeq E_Y$ over $G_{\mathbb Q}$.
\end{proposition}
\begin{proof}
We checked the entire proof of F149-ZCF4 and reran its finite
certificate after reading the program. The branch values of
$C\to\mathbb P^1$ are $0,1,-1,i,-i,\infty$. Its eight affine
singularity ideals are the unit ideal, all eighteen bilinear slice
determinants at a branch value are nonzero, and all 216 triples of
branch values give nonzero values. In local ramified coordinates
$x-b=z^2$ these checks respectively cover zero, one or two, and
three ramified coordinates, proving smoothness. The divisor is
ample, as the pullback of $\mathcal O(1,1,1)$ by a finite map.
The geometric Lefschetz theorem gives connectedness and an
isomorphism on geometric fundamental groups; the identity on
$G_{\mathbb Q}$ gives the arithmetic isomorphism. Finally
$\dim H=2<3=\dim Y$, which proves the failure of dominance through
the displayed actual morphism. The curve-target replacement is
Proposition~\ref{F-open-Hom}; no claim that it exhausts all
higher-dimensional replacements is made.
\end{proof}

\begin{remark}[F.6 --- Open: the remaining (F1) and field-reduction steps]
\label{F-open-main}
The unrestricted higher-dimensional sufficiency assertion in
F(b)(6) and the full reduction in F(c) from number fields remain
unproved. Proposition~\ref{F-local-dichotomy} proves the
infinite-index assertion in F(b)(5), using the local theorem
supplied in the prompt. Proposition~\ref{F-full-curve-reduction}
proves the exact universal curve-target (F1)--(F2) comparison.
The following records the obstacle in the earlier purely global
approach, now bypassed by the supplied local theorem. Its missing
group-theoretic step was this: for a continuous
$G_K$-map $f:E_V\to E_U$ satisfying (b), prove that the finitely
generated closed group $f(P_V)$ is either $1$ or open in $P_U$.
The elasticity theorem applies when this image is normal in an
\emph{open geometric subgroup}; normality in $f(E_V)$ does not
provide that hypothesis. Passing to the arithmetic normalizer leaves
precisely that unresolved input. The open Hom theorem cannot be
applied before openness has been supplied. The hyperplane example
above also prevents replacing this input by an unrestricted
higher-dimensional dominance assertion.

We checked Bresciani's published field-reduction paper
\cite{F-BrescianiFiniteType}: its proved range includes genus at most
two and curves with a suitable geometric finite cover mapping to a
constant abelian variety. Its general reduction retains an additional
Bombieri--Lang input. We have not proved that input or removed the
cover hypothesis. Spreading a section gives fibre section classes;
the unproved step in the unrestricted case is constructing a generic
rational point from that system. Accordingly the computations and
centralizer identities above supply no result labelled a proof of
this step. The first textual assertion in the 1983 notes not supplied
by our curve comparison is (F1) for arbitrary positive-dimensional
sources and elementary higher-dimensional targets, scan 48. This
identifies the scope not covered by that comparison; it is not a
claim that the entire 1981 manuscript lacks all higher-dimensional
ideas.
\end{remark}

\subsection{F(d): fixed points and descent from the cyclotomic kernel}

\begin{proposition}[F.7 --- Proved: fixed pro-points and sections]
\label{F-fixed-dictionary}
The completed geometric universal cover of a hyperbolic $U/K$
has at most one fixed geometric pro-point under a section restricted
to $H_K$. For a section over $G_K$ its $H_K$-fixed and $G_K$-fixed
sets coincide. An interior fixed pro-point gives precisely its
rational-point section, and a boundary fixed pro-point gives precisely
a section in its cusp packet. Thus the existence assertions in (F2)
and (F3), Conjecture B with its $H_K$ strengthening, have identical
content under these comparison maps.
\end{proposition}
\begin{proof}
The fibre of the universal cover over an interior geometric point
is a simply transitive $P_U$-set. Its arithmetic stabilizer maps
isomorphically to the point's Galois group. At a boundary pro-point
the stabilizer is $D_c=N_{E_U}(I_c)$. These two stabilizer
computations prove both maps of the dictionary. For uniqueness,
two pro-points differ at a finite covering level. Refine its
subgroup to a characteristic open subgroup with positive-genus
compactification. Such a cover exists: in genus zero pull back a
Fermat cubic along three distinct punctures; in positive genus use
the original curve. A further Galois and characteristic refinement
preserves positive genus by Riemann--Hurwitz. The given section
descends the cover over the cyclotomic field. The two fixed points
would give distinct rational points in its compactification with
the same section. The Jacobian argument in
Proposition~\ref{F-centralizer} excludes this. Normality of $H_K$
then makes a nonempty singleton fixed by all of $G_K$. This checks
CK D0a--D0b and uses no fixed-point existence assumption.
\end{proof}

\begin{proposition}[F.8 --- Refuted as literally stated: the discrete finite-group assertion]
\label{F-discrete}
The assertion in F(d) for \emph{every} finite group is false. The
identity splitting for $G=1$ fixes the entire universal cover.
For a nontrivial finite faithful group of orientation-preserving
hyperbolic isometries of the disk, the common fixed set is exactly
one interior point and contains no ideal boundary point.
Orientation-reversing lifts can instead fix a geodesic.
\end{proposition}
\begin{proof}
For $G=1$ the extension is the surface group itself and the unique
map $1\to\pi_1(U)$ is a splitting; its action is the identity on
every point. For the repaired assertion choose any finite orbit
in the disk and minimize the sum of its squared hyperbolic
distances. Strict convexity and properness give a unique minimum,
fixed by the group. Conjugate it to the centre. Every
orientation-preserving isometry fixing the centre is a rotation;
a nonidentity rotation fixes no other disk point and no point on
the ideal circle. This proves the unique fixed-point statement
for the specified action and also proves that such a finite group
is cyclic. For an orientation-reversing example, complex
conjugation $z\mapsto\overline z$ of the disk fixes its real
diameter. Take the real tripod with real point $2$; the lift of
complex conjugation fixing a chosen lift of $2$ is exactly a
reflection of its uniformizing disk and gives the corresponding
splitting by $C_2$. These examples specify the actual actions and
comparison maps. An arbitrary outer action requires a realization
before the isometry argument applies; finite-order and
orientation hypotheses cannot be omitted from that argument.
\end{proof}

\begin{proposition}[F.9 --- Proved: continuous B-prime actions acquire a quotient]
\label{F-continuous-action}
Let $\mathcal U$ be the geometric universal pro-cover of a hyperbolic
curve over $\overline K$. A semilinear $G_K$-action on $\mathcal U$
for which each function has open stabilizer normalizes the deck
group of a finite hyperbolic quotient defined over $K$. It is the
action of a continuous section for that quotient. The quotient can
be chosen to cover any specified initial level.
\end{proposition}
\begin{proof}
We checked F149-ZCF12, including descent of the open curve. Let
$\Omega$ be the union of covering function fields and let
$\overline K(f_1,\ldots,f_m)$ be an initial level. The orbits of
the $f_i$ are finite. The field $L$ they generate together is
therefore a finite extension of the initial field inside $\Omega$,
invariant under $G_K$. Its curve is a finite \'etale cover of the
initial level, and $\operatorname{Aut}_L(\Omega)$ is its deck group.
Normalization of this deck group descends every action map to
the quotient curve: factor it through a common finite Galois
cover, then descend the resulting deck-invariant map. Thus the
action preserves the actual open curve and its smooth proper
compactification with boundary. Choose a finite coefficient field
and intersect the stabilizers of finitely many defining generators;
its normal core is an open normal subgroup acting only on scalars.
The remaining finite semilinear action is an effective descent
datum. One explicit effectiveness proof embeds the marked pair
by $\omega(D)^{\otimes3}$ and descends its section space and
homogeneous ideal. The descended quotient is defined over $K$,
and the action on $\Omega$ fixes its function field. It is a
section of its arithmetic group, continuous by the stipulated
open stabilizers. Conversely every continuous section has those
stabilizers. Without this continuity requirement the finite-orbit
argument is unavailable. Proposition~\ref{F-discontinuous-action}
exhibits an abstract action for which no continuously descended
finite \'etale quotient exists, while its own fixed point still
exists. The universal abstract fixed-point assertion remains open.
\end{proof}

\begin{proposition}[F.10 --- Proved: the cyclotomic extension never splits virtually]
\label{F-no-cyclo-split}
For every number field $K$ and every open subgroup $Q'\subset Q_K$,
the extension
\[
 1\longrightarrow H_K\longrightarrow\chi_K^{-1}(Q')
       \longrightarrow Q'\longrightarrow1
\]
has no splitting, even as an abstract group extension.
\end{proposition}
\begin{proof}
The subgroup $Q'$ is open in $\widehat{\mathbb Z}^{\times}$, so it
contains a congruence kernel modulo some positive integer $m$.
Choose a prime $q\nmid m$ with $q\equiv1\pmod3$. There are such
primes outside every finite set: for a multiple $a$ of that set
and of $3$, any prime divisor $q$ of $a^2+a+1$ is outside the set,
is not $3$, and the order of $a$ modulo $q$ is $3$.
The factor $\mathbb Z_q^{\times}$ contains an element of order
$3$, obtained by the simple-root lift of an element of order $3$
in $\mathbb F_q^{\times}$. Put this in the $q$-coordinate and put
$1$ in every other coordinate. The resulting element of $Q'$
has order $3$. A splitting would embed it into $G_K$. The
Artin--Schreier theorem on finite subgroups of an absolute Galois
group in characteristic zero gives only the trivial group and
groups of order two. In particular $G_K$ has no element of order
three, a contradiction. This is the torsion argument requested
in the addition to F(d), with the open-subgroup quantifier included.
\end{proof}

\begin{proposition}[F.11 --- Proved: the exact obstruction to extending a cyclotomic section]
\label{F-descent-obstruction}
Let $1\to P\to E\xrightarrow{p}G_K\to1$ be a profinite
extension and $t:H_K\to E$ a continuous section over $H_K$.
Suppose its $P$-conjugacy class is $Q_K$-invariant. Put
\[
 C=C_P(t(H_K)),\qquad N=N_E(t(H_K)),\qquad W=N/t(H_K).
\]
There is an exact sequence
\[
 1\longrightarrow C\longrightarrow W\longrightarrow Q_K
       \longrightarrow1.                                      \tag{F.1}
\]
Extensions of the actual section $t$ to $G_K$ correspond bijectively
to continuous splittings of (F.1); conjugation by $C$ is respected.
In particular, $C=1$ gives a unique extension. For a cusp of a
hyperbolic curve, a choice of tangential origin identifies $C$ with
$T$, and the obstruction is the transgression
\[
 [o_t]\in H^2_{\rm cts}(Q_K,T),\qquad
 o_t(q,r)=-c\bigl(u(q)u(r)u(qr)^{-1}\bigr),                    \tag{F.2}
\]
where $t(h)=(c(h),h)$ in $T\rtimes G_K$, and $u$ is a continuous
set-theoretic transversal for $G_K\to Q_K$. Its vanishing is
necessary and sufficient; the extension classes form a torsor
under $H^1_{\rm cts}(Q_K,T)$ when nonempty.
\end{proposition}
\begin{proof}
For each $g\in G_K$, invariance supplies a lift in $E$ which
normalizes $t(H_K)$, by multiplying any lift by a suitable element
of $P$. Thus $N\to G_K$ is onto. Its kernel is $C$: an element of
$P$ normalizing $t(H_K)$ must centralize each $t(h)$, since
projection to $H_K$ distinguishes its elements. The subgroup
$t(H_K)$ is normal in $N$ and meets $C$ trivially. Quotienting
gives (F.1), as an exact sequence of compact groups.

A section extending $t$ lies in $N$ and descends to a splitting
of (F.1). Conversely, the inverse image in $N$ of the graph of a
splitting maps isomorphically to $G_K$: each fibre contains one
element, since its possible discrepancy lies in both $C$ and
$t(H_K)$. The compact-to-Hausdorff bijection is a homeomorphism.
Its inverse is the required section. This proves the criterion
without assuming a splitting of $G_K\to Q_K$.

For the cusp, the normalizer formula identifies $N$ with $D_c$;
invariance of $t$ is the equivariance of the homomorphism
$c:H_K\to T$. The multiplication of representatives
$(0,u(q))$ in $D_c/t(H_K)$ gives the error
$(0,u(q)u(r)u(qr)^{-1})$, whose image in the kernel is the negative
of $c$ of that element. This proves (F.2), including its sign.
Associativity proves the cocycle identity, and changing $u$ or
the representatives adds a continuous coboundary. The groups here
are countably based; a quotient of a countably based profinite
group by a closed subgroup admits a continuous transversal. To
construct one, choose successively finite clopen partitions of
the quotient subordinate to the images of cosets at level $n$,
refine the choices compatibly, and take the unique nested coset
limit. The image of every coset is open, which supplies each
finite refinement. Hence (F.2) is a continuous cocycle. Altering
the representatives by a continuous $1$-cochain kills it exactly
when the resulting representatives multiply as a section.
Two such choices differ by a $1$-cocycle, and $T$-conjugation
changes it by a coboundary. This proves the torsor assertion.
\end{proof}

\begin{proposition}[F.12 --- Proved: an invariant cuspidal class with nonzero descent obstruction]
\label{F-one-plus-i}
Let $F=\mathbb Q(\mu_\infty)$ and let $U=\mathbb P^1\setminus
\{0,1,\infty\}$. In the packet at $0$, the tangential class of
parameter $1+i$ over $F$ is invariant under
$Q=\operatorname{Gal}(F/\mathbb Q)$ and does not extend to a section
over $\mathbb Q$. Its class (F.2) is nonzero of order exactly two.
\end{proposition}
\begin{proof}
Every root of unity has roots of every order in $F$, so its Kummer
class in $\widehat{F^*}$ is zero. The two conjugates of $1+i$
have ratio $-i$, proving invariance. Let $c=\kappa_F(1+i)$.
The exact identity $(1+i)^2=2i$ gives
$2c=\operatorname{res}\kappa_{\mathbb Q}(2)$, and therefore kills
its transgression by two.

We compute the kernel of the restriction map, rather than infer
nonextension from a finite approximation. Inflation--restriction
identifies it with $H^1(Q,T)$. In fact
\[
 H^1(Q,T)=\mu(\mathbb Q)=\{1,-1\}.                            \tag{F.3}
\]
For completeness, apply the sequence
$0\to\mathbb Z_\ell(1)\to\mathbb Q_\ell(1)\to
\mathbb Q_\ell/\mathbb Z_\ell(1)\to0$.
The invariants of $\mathbb Q_\ell(1)$ vanish. Its first cohomology
also vanishes: for a continuous cocycle $b$ and a central $q$ with
$\chi_\ell(q)\ne1$, commutativity gives
$(\chi_\ell(q)-1)b(r)=(\chi_\ell(r)-1)b(q)$, explicitly a
coboundary after division in $\mathbb Q_\ell$. Thus the first
cohomology of $\mathbb Z_\ell(1)$ is the invariant torsion in
$\mathbb Q_\ell/\mathbb Z_\ell(1)$. The full cyclotomic image
fixes no such torsion for odd $\ell$, and fixes exactly its order-two
subgroup for $\ell=2$. Taking the product proves (F.3), whose
nontrivial class is the Kummer class of $-1$.

An extension of the class $c$ would have to be cuspidal at $0$:
it normalizes the nonzero centralizer $I_0$ of its restriction.
Its Kummer class $a\in\widehat{\mathbb Q^*}$ would satisfy
$\operatorname{res}(2a-\kappa_{\mathbb Q}(2))=0$.
By (F.3) the difference is a sign. The completed valuation at $2$
would then give $2v_2(a)=1$ in $\widehat{\mathbb Z}$, impossible
already modulo two. Consequently the obstruction is nonzero;
its order is two by the first paragraph. The attached exact
program verifies the two algebraic identities and the valuation
congruence obstruction. Its $D_8$ computation records the familiar
finite non-split central-extension phenomenon, but the proof of
nonextension above is the full profinite valuation argument.
\end{proof}

\subsection{F(e): function fields and the lifting obstruction}

\begin{proposition}[F.13 --- Proved: birational centralizers and the written formulas]
\label{F-birational}
Let $X/K$ be a smooth proper geometrically connected curve over a
finitely generated characteristic-zero field and let $L=K(X)$.
A section $s:G_K\to G_L$ has nontrivial centralizer in
$G_{\overline K(X)}$ of $s(H_K)$ precisely when it lies in the
decomposition group of a rational place. There the centralizer is
$I_v\simeq T$, both geometrically and arithmetically. Formulas
(32),(33) on scan 127 are false with their written $L$; for the
point section of a hyperbolic model $U$ containing that place,
their $U$ versions are true.
\end{proposition}
\begin{proof}
We checked CK B7--B8. The maps are the continuous surjections
$q_U:G_L\to\pi_1(U)$, as $U$ runs through the nonempty opens of
$X$; their inverse limit is $G_L$, since every finite separable
extension is unramified over some such open. A nonidentity
centralizer element survives in one hyperbolic open and every
smaller open. Proposition~\ref{F-centralizer} gives its unique
fixed cusp at each of these levels. Those cusps are compatible:
a cusp mapping to an interior point would kill its nonzero
inertia image. They determine compatible points on all finite
normalizations of $X$, hence a valuation of $\overline L$, fixed
by $s(G_K)$. Its point on $X$ is $K$-rational. Conversely a
decomposition section centralizes $I_v$ on $H_K$. The local
maps identify this same $I_v$ with its inertia image on every
open deleting the place; inverse limits give equality of the
centralizers. Projection to $G_K$ kills the arithmetic centralizer
by the arithmetic calculation used above. A model containing the
place kills $I_v$ and gives a point section, whose centralizer is
zero by Proposition~\ref{F-centralizer}.

Explicitly take $L=\mathbb Q(t)$, the $t$-adic place $v$, and
$U=\mathbb P^1\setminus\{1,2,\infty\}$. Fixing compatible roots
of $t$ gives a decomposition section. Its centralizer contains the
inertia detected by $t^{1/2}$ and equals $T$; its image on $U$ is
the section of $0$. The comparison is exactly $q_U$, which kills
that ramification. Visual inspection of scan 127 supports the
subscript $L$ in both formulas. The mathematical counterexample
does not authorize silently changing that reading.
\end{proof}

\begin{proposition}[F.14 --- Known, with exact quantifiers: universal cuspidalization]
\label{F-cuspidalization}
For a hyperbolic curve $X/k$, $k/\mathbb Q$ finitely generated,
the assertion that all sections after every finitely generated
extension of $k$ are geometric or cuspidal is equivalent to the
assertion that all those sections lift through the function-field
Galois group. The references arXiv:2108.13397 and Inventiones
235 (2024), 129--150 name the same paper.
\end{proposition}
\begin{proof}
We opened the publisher's version of Bresciani and checked
Theorems A and B~\cite{F-BrescianiBirational}. For a section $s$
over a finitely generated $K/k$, its base change to $K(t)$ is
the actual map
\[
 g\longmapsto(s(\operatorname{res}g),g)
 \quad\hbox{in}\quad
 \pi_1(X_K)\times_{G_K}G_{K(t)}=\pi_1(X_{K(t)}).
\]
The universal lifting assertion applies to this section. Theorem A
then gives geometricity or cuspidality of $s$. Conversely a point
section lifts through its rational-place decomposition sequence,
and a cusp section already specifies such a local splitting.
Geometric conjugators lift through the surjective generic-point
map. This proves the equivalence on classes and on chosen
representatives. The additional $K(t)$ base change is present in
Theorem A and is indispensable in this reduction; the theorem
does not establish the same conclusion from ordinary birational
liftability alone.
\end{proof}

\begin{remark}[F.15 --- Open: the remaining valuative and higher-dimensional lifting claims]
\label{F-open-valuations}
For the tripod over $\mathbb Q$, the exact lifting problem still
left by our work is the splitting of
\[
 1\longrightarrow\ker(G_{\mathbb Q(t)}\to E_U)
 \longrightarrow G_{\mathbb Q(t)}\times_{E_U,s}G_{\mathbb Q}
 \longrightarrow G_{\mathbb Q}\longrightarrow1
\]
for an arbitrary section $s$. All geometric and cuspidal sections
have such a lift by the local construction, but the inertia
description of its kernel gives no complement for an arbitrary
$s$. The centralizer test in Proposition~\ref{F-birational}
recognizes a decomposition group once nonzero centralizer has
been supplied; it does not supply that nonzero element.

The full-rank structure and the exact normalizer are established
in Proposition~\ref{F-valuation-structure}. The finite-residue
qualification needed to produce germs, and its failure for an
infinite algebraic residue extension, are established in
Proposition~\ref{F-valuation-germs}. Restricting a decomposition
section along $K\subset L_1\subset L$ sends its valuation to
$V\cap L_1$ and proves the forward compatibility on scans
125--126. Reversing it would require a single compatible invariant
valuation above all the one-variable restrictions, and no such
construction is supplied. In particular, selecting a place for
each subfield does not itself show that the resulting valuation
rings are compatible under inclusion.

For F(e)(4), a section of a polycurve gives a section on its base.
A birational lift of that base section has domain $G_K$, whereas
the curve-fibre lifting assertion over $K(B)$ has domain
$G_{K(B)}$. The morphisms displayed here do not extend that
domain. Applying the latter assertion without this extension
would leave precisely the missing step identified in the prompt.
Proposition~\ref{F-proper-polycurve-lifting} supplies the
implication for towers proper except possibly at their last stage,
using actual point fibres and full-flag valuations. The case with
a cuspidal projected section on an affine intermediate base remains
unproved in this return.
\end{remark}

\subsection{F(f): audit of the potential-good-reduction criterion}

\begin{proposition}[F.16 --- Known, with checked comparison: condition (e)]
\label{F-reduction}
Let $R$ be a henselian discrete valuation ring whose residue field
$k$ is finitely generated over $\mathbb Q$, and let $U/K$ be a
hyperbolic curve over its fraction field. The residue action on
its full geometric profinite group obtained from any splitting
$G_k\to G_K$ is realized by a smooth hyperbolic curve over $k$
precisely when $U$ has potentially good marked reduction. Allowing
a finite residue extension gives the same condition. This is the
corrected condition (e) of cote 149, scans 100--103.
\end{proposition}
\begin{proof}
We checked the complete proofs of F149-F17.6 and F17.7. The
essential comparison is between the full group of the generic
curve and the logarithmic special fibre, rather than just their
homology. Illusie, Theorem 7.6~\cite{F-Illusie}, identifies covers
of the generic interior with finite Kummer \'etale covers of the
log-regular model; all boundary ramification is tame in residue
characteristic zero. Hoshi, Corollary 3.9~\cite{F-HoshiLog},
identifies the latter covers with those of the strict log special
fibre after completion. Completion induces the same residue
quotient and the same inertia $T$, hence the same absolute Galois
group here. The curve geometric groups also remain unchanged in
characteristic zero. These are compatible isomorphisms over the
base log Galois group.

At a split node of thickness $m$, the monoid chart is
$X+Y=mT_0$. Its arithmetic local group has geometric torus
\[
 \{(a,b,c)\in T^3:a+b=mc\},
\]
and projection to base inertia is $(a,b,c)\mapsto c$. The kernel
is $(a,-a,0)$, an injected node inertia. Pulling back along any
specified residue splitting preserves surjectivity to $G_k$ and
the cyclotomic action on this kernel. Mochizuki, Example 2.5 and
Proposition 1.2~\cite{F-MochizukiCombinatorial}, together with the
node injection in Hoshi--Mochizuki Lemma 2.11~\cite{F-HoshiMochizukiNodes},
identifies these as distinct maximal procyclic node groups.
Thus every node contributes a cyclotomically normalized class
in addition to the cusp classes.

We verify the numerical exclusion used to compare with an
arbitrary smooth realizing curve. After a finite residue extension
fixing the dual graph, a cycle in that graph gives a nonzero
trivial quotient of $H_1(U_{\overline K},\mathbb Q_\ell)$.
A smooth curve has none: its toric homology has character
$\chi_\ell$, while its Jacobian Tate module and its dual have
zero cyclotomic-kernel invariants by the arithmetic argument of
Proposition~\ref{F-centralizer}. For a tree with $r>0$ cusps,
eliminating end vertices from the boundary relations gives
\[
 0\to\mathbb Q_\ell(1)^{r-1}\to H_1(U_{\overline K},\mathbb Q_\ell)
 \to\bigoplus_v V_\ell\operatorname{Jac}(C_v)\to0.
\]
The dimension of maps from $\mathbb Q_\ell(1)$ is exactly $r-1$.
A smooth realizing curve would therefore have exactly $r$ cusps,
but Stix's cyclotomic normalizer criterion gives only $r$ such
classes, contradicting the $r+n$ classes from $n>0$ nodes. For
a proper curve, its nonzero $H^2(-,\mathbb Z/\ell)$ first forces
the realizing curve to be proper, and a single node class already
contradicts the same criterion. This proves nonrealization after
every finite residue extension for stable bad reduction.

For potentially good reduction choose a finite Galois extension
$K_1/K$ giving a smooth stable marked model. Uniqueness of that
model extends the Galois action to it and gives descent maps on
its special pair over $k_1$. Compose the specified splitting with
$G_K\to\operatorname{Gal}(K_1/K)$ and use these maps to twist
the pair. After a finite extension of $k$ this is an ordinary
finite Galois descent datum, effective by the marked
$\omega(D)^{\otimes3}$ embedding. Its descended curve over $k$
has the required action: the generic action factors through
\[
 \operatorname{Gal}(K_1/K)
       \times_{\operatorname{Gal}(k_1/k)}G_k,
\]
and restriction along the splitting is exactly the descent action.
Conversely, choose stable reduction after a finite field
extension, intersect the splitting's inverse image with its open
Galois subgroup, and apply the stable-bad exclusion above. Nodes
survive further stable base change. Hence realization implies a
smooth stable special fibre. This proves both directions with the
specified full-profinite comparison, and verifies the claim being
used from F149 without changing its quantifiers.
\end{proof}

\subsection{F(g): geometric tangents, local conditions, and ideles}

\begin{proposition}[F.17 --- Proved: the two torsors and their coordinate ambiguity]
\label{F-torsors}
At a rational cusp over a finitely generated field $K/\mathbb Q$, the packet is a torsor
under $\widehat{K^*}=H^1(K,T)$, and the geometric tangent vectors
give its injective $K^*$-subtorsor. For a smooth connected separated
finite-type $\mathbb Q$-scheme $S$ and an injected normal-line
extension as in scan 85, its class is the Kummer first Chern class
of the normal line $\mathcal N$; it splits precisely when
$\mathcal N$ is trivial. In that event the geometric frame torsor
under $\mathcal O(S)^*$ embeds in the splitting torsor under
$H^1(S,T)=\widehat{\mathcal O(S)^*}$.
\end{proposition}
\begin{proof}
For a local parameter $z$, compatible roots of $z$ give an origin
in the packet. Fix the convention
$\zeta:z^{1/n}\mapsto\zeta_nz^{1/n}$ and
$\kappa(a)(\sigma)=\sigma(a^{1/n})/a^{1/n}$.
The section fixing roots of $az$ differs from the section fixing
roots of $z$ by $-\kappa(a)$. A tangent vector $v$ with
$z(v)=a$ instead has coordinate $+\kappa(a)$. These explicit
inverse maps identify the conormal-frame and tangent conventions.
The Kummer kernel in
a finitely generated characteristic-zero field is trivial:
divisibility by all integers kills every divisorial valuation,
and then reduces to its number-field constant units, where it
again forces the identity. This proves injectivity.

For $S$ choose local normal frames $e_j=g_{ij}e_i$. Lifts of these
transition functions through their compatible roots have the
triple-overlap defect defining $c_1(\mathcal N)$; these are exactly
the composition defects of the local normal-tube splittings.
This is the sign convention of F149-C6. The conormal line has
the negative class. A splitting kills this class. The Kummer
sequence injects $\operatorname{Pic}(S)/n$ into
$H^2(S,\mu_n)$, so $\mathcal N$ is divisible by every $n$ in
$\operatorname{Pic}(S)$. Here $\operatorname{Pic}(S)$ is finitely
generated: choose a smooth proper compactification after
resolution, use finite generation of its rational Picard group,
and quotient by the finitely many boundary divisors. Thus
$\mathcal N$ is trivial. The converse is given by a global frame.
Taking inverse limits in the degree-one Kummer sequence gives
$H^1(S,T)=\widehat{\mathcal O(S)^*}$ because
$T\operatorname{Pic}(S)=0$; the finite root-of-unity invariants
give no derived-limit obstruction in degree one. Embedding
$\mathcal O(S)^*$ into the function field gives its injectivity
by the same Kummer-kernel argument. These are the checked steps
of F149-C6--C8, with left injection explicitly retained as part
of the normal-line extension's data.
\end{proof}

\begin{proposition}[F.18 --- Proved: the abstract local extension does not select its geometric subtorsor]
\label{F-twist}
Automorphisms of $T\rtimes G_K$ inducing the identity on both
$T$ and $G_K$ act transitively on its section classes. Within
$E_U$, the cusp packet embeds in the full section set, so a
chosen parameter gives a uniquely determined completed Kummer
coordinate for each cuspidal class. Without that choice its
coordinate is defined up to translation by $K^*$.
\end{proposition}
\begin{proof}
For a continuous cocycle $b:G_K\to T$ the formula
$(a,g)\mapsto(a+b(g),g)$ is a continuous group automorphism;
the cocycle identity is exactly multiplicativity. It carries the
zero section to the class of $b$, proving transitivity. Since
$K^*$ is a proper subset of $\widehat{K^*}$, that automorphism
group cannot preserve the proposed proper subtorsor. The packet
injection inside $E_U$ follows directly from
Proposition~\ref{F-centralizer}: a geometric conjugator between
two sections at the same lifted cusp must normalize $I_c$,
hence lies in $I_c$. Therefore precisely the local coboundary
ambiguity remains. Replacing the chosen parameter multiplies it
by a nonzero scalar in its first-order conormal line and shifts
the coordinate by that scalar's Kummer class.
\end{proof}

\begin{proposition}[F.19 --- Proved: all everywhere locally geometric global Kummer classes]
\label{F-local-global-units}
For a number field $K$ put $E=\mathcal O_K^*$ and define
$\mathcal L_K\subset\widehat{K^*}$ by the two requirements:
its image in $\widehat{K_v^*}$ belongs to $K_v^*$ at every finite
place $v$, and it belongs to $\mathcal O_v^*$ at almost every
finite place. Then
\[
 \mathcal L_K=K^*\,\widehat E,\qquad
 \mathcal L_K/K^*\simeq\widehat E/E.                          \tag{F.4}
\]
For $K=\mathbb Q$ this is exactly $\mathcal L_{\mathbb Q}=\mathbb Q^*$.
For every infinite-order unit $\varepsilon$ and
$\alpha\in\widehat{\mathbb Z}\setminus\mathbb Z$, the class
$\varepsilon^\alpha$ is in $\mathcal L_K\setminus K^*$.
\end{proposition}
\begin{proof}
At a finite place, a uniformizer gives
$\widehat{K_v^*}=\widehat{\mathbb Z}\times\mathcal O_v^*$.
The unit factor is already complete because it is profinite
abelian: its map to the inverse limit of power quotients is
injective by its finite quotients and surjective by compactness.
Thus membership in $K_v^*$ means exactly that the valuation is
an ordinary integer; the almost-all-unit condition means these
integers have finite support.

Let $D=\bigoplus_v\mathbb Z[v]$ be the finite divisor group and
$P\subset D$ its principal divisors. The valuation sequence
$1\to E\to K^*\to P\to0$ splits as an abstract abelian sequence,
since $P$ is free abelian. Consequently
$\widehat{K^*}=\widehat E\times\widehat P$ after such a choice.
The inclusion $P\to D$ induces an injection
$\widehat P\to\widehat D$. Indeed $D/P=\operatorname{Cl}(K)$
has a finite exponent $h$, and
$(hn)D\cap P\subset nP$; equivalently the induced power
topologies are cofinal. The same conclusion follows by choosing
a basis for the finite-index subgroup $P$ of the free abelian
group $D$. An integral finite divisor in the image of
$\widehat P$ is principal: its image in the finite class group
is zero, because that class-group map factors through $D/hD$.
Hence a class $a\in\mathcal L_K$ has the divisor of some
$b\in K^*$. Its quotient by $b$ lies in $\widehat E$, proving
the first equality. Conversely the image of $\widehat E$ at
every finite place lies in the compact group $\mathcal O_v^*$,
so every element of $K^*\widehat E$ has the required properties.
The intersection $K^*\cap\widehat E$ is exactly $E$, giving
the quotient statement.

The finite generation $E\simeq\mu(K)\times\mathbb Z^r$
makes the last assertion exact. An infinite-order unit has a
nonzero vector $m\in\mathbb Z^r$. The vector $\alpha m$
cannot be integral unless $\alpha\in\mathbb Z$: a nonzero
coordinate gives $m_i\alpha=n_i$, and an element of
$\widehat{\mathbb Z}$ equal to a rational number is an integer.
This proves nonglobality while the first equality proves all
local conditions. For $\mathbb Q$, $E=\{1,-1\}$ is already
complete, so no extra classes occur.
\end{proof}

\begin{proposition}[F.20 --- Refuted and replaced: powers of one unit do not exhaust the defect]
\label{F-one-unit-false}
The added assertion in F(g)(3 bis) that all nonglobal everywhere
locally geometric classes are exactly the $\varepsilon^\alpha$
for one fixed infinite-order unit is false. The exact complete
replacement is (F.4). The associated idele classes of such locally
geometric classes, with matching real signs, belong to the
connected component of the idele class group.
\end{proposition}
\begin{proof}
In $K=\mathbb Q(\sqrt2,\sqrt3)$ take
$\varepsilon=1+\sqrt2$, $\eta=2+\sqrt3$, and let
$e_2\in\widehat{\mathbb Z}$ have $2$-component $1$ and every
other prime component zero. The class $\eta^{e_2}$ is locally
geometric and nonglobal by Proposition~\ref{F-local-global-units}.
It cannot be $\varepsilon^\alpha$: the norm to
$\mathbb Q(\sqrt3)$ takes $\varepsilon$ to $-1$ and $\eta$
to $\eta^2$. Thus such an equality would make
$\eta^{2e_2}$ torsion in the completed unit group, whereas its
nonzero $2$-adic unit-lattice vector makes it infinite order.
The exact unit norms are checked by the attached program.
Even in unit rank one, the actual set in (F.4) also contains
arbitrary $K^*$ multiples; its quotient is the one described
there, with all unit directions retained.

To check the connected-component assertion without assuming a
chosen reciprocity convention, approximate a completed unit by
ordinary units simultaneously in every finite-place compact unit
group and in its finite real-sign group. The corresponding idele
classes are represented by archimedean positive components alone,
after division by the principal unit idele. Those classes lie in
the image of the connected identity component of the archimedean
idele group. Passing to the limit remains in its closed connected
component. Multiplication by a principal $K^*$ element does not
change the idele class. Different positive archimedean magnitudes
also remain in that same component. This proves exactly the
asserted membership; it does not assign an unchosen canonical
archimedean representative to a completed Kummer class.
\end{proof}

\begin{proposition}[F.21 --- Proved with the necessary restriction: rational idelic gluing]
\label{F-ideles}
There is a canonical decomposition, using the real sign,
\[
 \varprojlim_n\mathbb A_{\mathbb Q}^{\times}/
                   (\mathbb A_{\mathbb Q}^{\times})^n
 \simeq\widehat{\mathbb Q^*}\times\widehat{\mathbb Z}^{\times}.
                                                                  \tag{F.5}
\]
A family of local completed Kummer classes belongs to its left
side precisely when, for every $n$, only finitely many finite
places have nonzero valuation modulo $n$. For such a family
$(b_v)$, let $a\in\widehat{\mathbb Q^*}$ have its real sign
and its completed valuations. Its unique second coordinate is
\[
 u_p=b_p\,\operatorname{res}_p(a)^{-1}\in\mathbb Z_p^*,
 \qquad u=(u_p)_p\in\widehat{\mathbb Z}^{\times}.             \tag{F.6}
\]
For local cuspidal sections at one fixed rational cusp with a
common rational tangent origin, $u=1$ is necessary and sufficient
for the family to be the restrictions of one global cuspidal
section.
\end{proposition}
\begin{proof}
For an actual idele $b$, put
$q=\operatorname{sgn}(b_\infty)\prod_p p^{v_p(b_p)}$.
There are finitely many nonzero valuations. Then
$b=q\,u\,r$ uniquely, with $u\in\prod_p\mathbb Z_p^*$
and $r\in\mathbb R_{>0}$. This gives the topological product
$\mathbb Q^*\times\widehat{\mathbb Z}^{\times}\times
\mathbb R_{>0}$. The last factor has all roots, and the compact
middle factor is complete for power quotients by the compactness
argument in Proposition~\ref{F-local-global-units}. Taking
inverse limits proves (F.5).

At each $n$, an idele class modulo $n$ has finitely supported
valuations modulo $n$, and every such family of local classes
has an idele representative: at the other places choose unit
representatives. A family locally an $n$th power is globally an
idele $n$th power by choosing roots at all places, which are
units almost everywhere. Thus these local maps are injective
and their stated images are exact. Passing to the inverse limit
proves the criterion and (F.6). Equivalently the positive prime
exponents of $a$ form a family in $\widehat{\mathbb Z}$ tending
to zero modulo every $n$, so their formal product is an element
of $\widehat{\mathbb Q^*}$. The remaining local valuations are
zero, giving the indicated unit factors.

The cusp packet maps are the Kummer restriction maps after the
common tangent origin is chosen. A global class therefore has
the coordinate $(a,1)$, and every such coordinate defines its
global cuspidal section. The packet injection in
Proposition~\ref{F-twist} proves necessity and uniqueness also
after passage to full arithmetic section classes. Changing the
common tangent origin by $q\in\mathbb Q^*$ changes $a$ by $q$
and leaves $u$ unchanged. Identifying this unit group with
$\operatorname{Gal}(\mathbb Q(\mu_\infty)/\mathbb Q)$ uses
the cyclotomic isomorphism; an Artin-map convention can invert
that isomorphism without affecting the obstruction $u=1$.
\end{proof}

\begin{proposition}[F.22 --- Refuted: arbitrary compatible local families have the idelic decomposition]
\label{F-ideles-unrestricted}
Compatibility of local Kummer classes modulo every $n$, by itself,
does not imply the decomposition in (F.5). Omitting the real
place also removes uniqueness of the two factors.
\end{proposition}
\begin{proof}
Take at each finite prime $p$ the actual local parameter $b_p=p$,
and real sign $+1$. These give compatible local section classes
at the rational cusp $0$ for every modulus. Their valuations are
$1$ at every prime, so already modulo two they have infinite
support. Proposition~\ref{F-ideles} excludes them from the
completed idele group. This is a full counterexample to the
unrestricted family clause in the added task. For finite places
alone the pairs $(a,u)$ and $(-a,-u)$ give identical local
products: the missing real sign had distinguished these pairs.
The all-place formulation of Proposition~\ref{F-ideles} and its
restrictedness condition are the exact replacements proved here.
\end{proof}

\begin{proposition}[F.23 --- Proved: geometric tangents are uniformizer specializations]
\label{F-specialization}
Let $K/\mathbb Q$ be finitely generated, let a smooth marked
curve over $K[[\pi]]$ have a rational cusp section, and let a
generic rational point meet this cusp with multiplicity $d>0$.
Write a parameter along the cusp as
\[
 z=a\pi^d(1+h),\qquad a\in K^*,\quad h\in\pi K[[\pi]].
                                                                  \tag{F.7}
\]
Specialization using the uniformizer splitting of the base
Galois group gives exactly the tangent packet class $\kappa_K(a)$.
Replacing $\pi$ by $b\pi(1+h')$ replaces it by
$\kappa_K(ab^{-d})$. Every geometric tangent class occurs in
this way. Thus in this precise smooth marked trait situation the
geometric subtorsor is exactly the set of such specializations.
\end{proposition}
\begin{proof}
The base splitting fixes compatible roots of $\pi$ and acts on
coefficients by $G_K$. The factor $1+h$ has a unique $n$th
root congruent to $1$ for every $n$, by Hensel's lemma in residue
characteristic zero; these roots are compatible. Pulling back the
Kummer cover $z^{1/n}$ along the point of (F.7) therefore leaves
precisely the cocycle from $a^{1/n}$. The inertia homomorphism is
multiplication by $d$, while the residue section coordinate is
$\kappa_K(a)$. This proves the formula at every finite level and
hence in $T$, and substitution of the new uniformizer proves
its stated change. Conversely the formal smooth parameter chart
at the cusp admits an arc with $z=a\pi$. Its generic point is
in the open curve and its specialization is the desired class.
For the tripod this is the explicit point $x=a\pi$; formal
smoothness gives the same construction for any smooth rational
cusp. This is the exact comparison asked in scans 85--95.
\end{proof}

\begin{remark}[F.24 --- Open: a formula using only GT automorphisms]
\label{F-open-tangents}
The local conditions do characterize the geometric tangent
subtorsor over $\mathbb Q$ by
Proposition~\ref{F-local-global-units}. Over a general number
field their exact residual defect is $\widehat E/E$.
The cocycle-twist automorphisms in Proposition~\ref{F-twist}
prove why the isolated local extension does not remove this
defect. We have not produced or refuted a single
Grothendieck--Teichm\"uller relation selecting every prescribed
parameter $t\in\mathbb Q^*$ as demanded by F(g)(5). The
specialization characterization in
Proposition~\ref{F-specialization} is proved with the indicated
family and uniformizer data; it does not erase those data to
produce an additional intrinsic equation in the automorphism
group of the single tripod.
Theorem~\ref{F-two-Hom-tangents} below does give an intrinsic
global relation inside that arithmetic group when two open
homomorphisms from one of its open subgroups are allowed.
It recognizes the entire ordinary tangential subtorsor over
every number field, including the unit-completion defect above.
The remaining question in this remark concerns the expressly
more restrictive automorphism-only form of the relation.
\end{remark}

\subsection{Further exact comparisons and remaining reductions}

\begin{proposition}[F.25 --- Known and proved: inertia and its normalizer in every rank]
\label{F-valuation-structure}
Let $L/K$ be finitely generated of characteristic zero, and let
$v$ be a valuation of $L$ trivial on $K$. Fix its prolongation to
$\overline L$, and write $\Gamma_v=v(L^*)$. There are a split
continuous exact sequence and canonical character identifications
\[
 1\longrightarrow I_v\longrightarrow D_v
   \longrightarrow G_{\kappa(v)}\longrightarrow1,
 \qquad
 I_v=\operatorname{Hom}((\Gamma_v\otimes\mathbb Q)/\Gamma_v,
                         \mu_\infty)
     =\operatorname{Hom}(\Gamma_v,T).
\]
The splitting has no preferred choice, and the action on $I_v$
is cyclotomic. Noncanonically,
\[
 I_v\simeq\prod_\ell\mathbb Z_\ell(1)^{r_\ell},\qquad
 r_\ell=\dim_{\mathbb F_\ell}\Gamma_v/\ell\Gamma_v
 \leq\operatorname{rank}_{\mathbb Q}\Gamma_v
 \leq\operatorname{trdeg}_K L-\operatorname{trdeg}_K\kappa(v).
\]
For $I_v\ne1$, let $\Delta$ be the largest divisible convex
subgroup of $\Gamma_v$ and let $w=v/\Delta$ be the coarsening.
Then, as actual subgroups for the specified prolongations,
\[
 I_w=I_v,\qquad N_{G_L}(I_v)=D_w.
\]
Thus $N_{G_L}(I_v)=D_v$ exactly when the residual valuation
$v/w$ on $\kappa(w)$ is henselian. In particular it holds for
rank one with nonzero inertia. For $I_v=1$ the normalizer is
all of $G_L$. These are the precise replacements for the
codimension and normalizer assertions on cote 149, scans
119--124 and 127.
\end{proposition}
\begin{proof}
We checked F149-D4--D5 and D10 against the named sources.
Over the strict henselization the residue field is algebraically
closed and every unit has all roots, by Hensel's lemma. Residue
characteristic zero makes every finite extension tame and
defectless; this is Kuhlmann, Theorem 1.2 and Corollary
2.9~\cite{F-KuhlmannTame}. Temkin, Theorem 3.4.4, describes
the tame extensions by the finite enlargements of the value
group~\cite{F-Temkin}. Explicitly the pairing is
\[
 (\sigma,v(a)/n)\longmapsto \sigma(a^{1/n})/a^{1/n}.
\]
It is independent of the representative because the quotient
of two representatives of the same value is a unit with an
$n$th root. This gives the first character group. Its image in
$\operatorname{Hom}(\Gamma_v,T)$ sends $\gamma$ to the
compatible sequence whose $n$th term evaluates the character
on $\gamma/n\bmod\Gamma_v$. This gives the second displayed
isomorphism, rather than an unidentified abstract equivalence.
The torsion group $(\Gamma_v\otimes\mathbb Q)/\Gamma_v$ is
divisible and has $\ell$-torsion $\Gamma_v/\ell\Gamma_v$.
Choosing compatible roots of a basis gives its sum of Pr\"ufer
groups and the stated dual product. Independence modulo $\ell$
implies rational independence, by clearing denominators and
canceling a largest common power of $\ell$. Finally, elements
with rationally independent values, together with units whose
residues form a transcendence basis, are algebraically
independent: in any proposed polynomial relation the monomials
of least value have a nonzero residue polynomial. This proves
the two rank bounds.

For completeness the splitting can be constructed in the chosen
algebraic closure. Choose a maximal algebraic extension $M$ of
the henselization with unchanged residue field. Its value group
is divisible: adjoining an $\ell$th root of an element with value
outside $\ell\Gamma_M$ has ramification degree $\ell$ and
unchanged residue, contradicting maximality. The maximal
unramified extension of $M$ has divisible value group and
algebraically closed residue. Defectlessness makes it algebraically
closed. Hence $G_M\to G_{\kappa(v)}$ is an isomorphism, and
its inclusion in $D_v$ supplies the splitting. This also checks
the characteristic-zero specialization of Haran--Jarden--Pop,
Proposition 7.3(a)~\cite{F-HaranJardenPop}. Conjugation in the
Kummer pairing acts on its roots of unity by the cyclotomic
character and leaves the value group fixed.

Convex subgroups form a chain, and the union of its divisible
members is divisible. For a convex subgroup $\Delta'$ the
coarsening inertia is the character subgroup obtained by pullback
from $\Gamma_v/\Delta'$. It equals $I_v$ precisely when
$\Delta'$ is divisible. Put $F=\overline L^{I_v}$. Its induced
valuation is henselian, with value group $\Gamma_v$ and
algebraically closed residue field. The coarsenings with
algebraically closed residue are precisely those with divisible
convex kernel, by the already proved inertia computation applied
to their residual refinements. Thus $w|F$ is the coarsest of
these coarsenings. It is coarsest among all henselian valuations
of $F$ with separably closed residue: two incomparable henselian
valuations have a common coarsening with separably closed
residue, by the F.~K.~Schmidt theorem in Kuhlmann, Theorem 9.44
and Corollary 9.45~\cite{F-KuhlmannHensel}. Such a common
coarsening would contradict the choice of $w$. Comparability
then forces every such valuation to refine $w$.

Every element normalizing $I_v$ acts on $F$ and preserves this
intrinsically specified coarsest valuation. Its unique
prolongation to $\overline L$ is therefore preserved, placing
the element in $D_w$. Conversely $D_w$ normalizes $I_w=I_v$.
The inverse image of the residual decomposition group of $v/w$
in $D_w$ is $D_v$, so equality is precisely henselianity of the
residual valuation. In rank one the convex subgroups are $0$
and $\Gamma_v$; nonzero inertia excludes divisibility of the
latter. This proves the last assertions and the actual comparison
of all normalizers.
\end{proof}

\begin{proposition}[F.26 --- Proved, with counterexample: the finite-residue requirement for germs]
\label{F-valuation-germs}
For a valuation $v$ as in Proposition~\ref{F-valuation-structure}
whose residue field is algebraic over $K$, the decomposition group
$D_v$ contains a section over an open subgroup of $G_K$ if and
only if $[\kappa(v):K]<\infty$. It contains a section over all
of $G_K$ precisely when $\kappa(v)=K$. Algebraicity alone
does not suffice, even for a DVR of $\mathbb Q(x,y)$.
With the finite-residue qualification, restriction to any
one-variable subfield is exactly the construction on scans
125--126: its valuation ring is $V\cap L_1$, and its section
lands in the corresponding rational-place decomposition group
after the finite residue extension.
\end{proposition}
\begin{proof}
Use the embedding $\kappa(v)\subset\overline K$ fixed by the
prolongation. The image of $D_v\to G_K$ is exactly
$G_{\kappa(v)}$. It contains an open subgroup precisely when
this residue extension is finite. The splitting in
Proposition~\ref{F-valuation-structure} then supplies the section
over that subgroup. The same argument with the full image proves
the assertion over $G_K$.

Here is an explicit infinite-residue example. Choose distinct
rational primes $p_1,p_2,\ldots$, put $a_j=\sqrt{p_j}$,
$F_m=\mathbb Q(a_1,\ldots,a_m)$, and $F_\infty=\bigcup_mF_m$.
Inside $F_\infty((t))$ set
\[
 x=t,\qquad y=\sum_{j\geq1}a_jt^j,
 \qquad L=\mathbb Q(x,y),
\]
and restrict the $t$-adic valuation. Its value group is
$\mathbb Z$ because $x=t$, and its residue field is contained
in $F_\infty$. The residue of $y/x$ is $a_1$. For $m\geq1$
form the rational norm polynomial
\[
 P_m(X,Y)=\prod_{\sigma\in\operatorname{Gal}(F_m/\mathbb Q)}
 \left(Y-\sum_{j=1}^m\sigma(a_j)X^j\right)\in\mathbb Q[X,Y].
\]
The identity factor at $(t,y)$ starts with $a_{m+1}t^{m+1}$.
For a nonidentity sign vector whose first changed index is $j$,
the first term is $2a_jt^j$; there are $2^{m-j}$ such vectors.
Consequently
\[
 v(P_m(x,y))=2^{m+1}-1,
 \quad
 \overline{x^{-(2^{m+1}-1)}P_m(x,y)}
 =2^{2^m-1}\!\left(\prod_{j<m}p_j^{2^{m-j-1}}\right)
   a_m a_{m+1}.
\]
Induction gives $F_m\subset\kappa(v)$ for every $m$, hence
$\kappa(v)=F_\infty$. Were $y$ algebraic over $\mathbb Q(t)$,
the fundamental inequality for the resulting finite valued-field
extension would bound its residue degree. The displayed residue
field has degrees $2^m$ for all $m$, a contradiction. Thus this
$L$ is the rational function field in two variables, and the
valuation is a DVR with infinite algebraic residue. Its inertia
is $T\ne1$, its normalizer is $D_v$, and neither group supplies
the asserted germ. The exact program
\texttt{infinite\_residue\_checks.py} verifies the first four
norm polynomials and their leading coefficients; the induction
above proves the infinite assertion.

For a subfield $L_1$, either $b$ or $b^{-1}$ belongs to
$V\cap L_1$ for every $b\in L_1^*$, and the same alternatives
express $b$ as a quotient of two elements of that ring. It is
therefore a valuation ring with fraction field $L_1$, and its
residue injects into $\kappa(v)$. Compatible prolongations give
the actual restriction maps $D_v\to D_{v|L_1}$ and
$I_v\to I_{v|L_1}$. If $L_1/K$ has transcendence degree one,
the restricted valuation is nontrivial, since its residue embeds
in an algebraic extension of $K$. Its center on the smooth proper
curve is a closed point and its ring is the DVR there. The
section therefore restricts to exactly that decomposition group.
This proves the manuscript's forward construction. On scan 126
we visually read ``fini'' after the struck-out ``alg''; that
qualification is mathematically necessary. The supplied first-pass
transcription reports the same insertion on scan 124, whose image
is not included in the supplied selection.
\end{proof}

\begin{proposition}[F.27 --- Proved: the one-variable equivalence has both lifting directions]
\label{F-birational-equivalence}
Fix a smooth proper geometrically connected curve $X/K$, with
$K/\mathbb Q$ finitely generated and $L=K(X)$. The conjunction
of (F4)(a) for $L K'/K'$ and (F4)(b) for every hyperbolic open
$U\subset X_{K'}$, for all finite $K'/K$, is equivalent to
(F2) for all those open curves. In this comparison (a) means
that every birational section lies in a rational-place
decomposition group; in transcendence degree one the place
and its fixed prolongation are unique. Proposition~\ref{F-birational}
identifies this (a), section by section, with nontrivial
cyclotomic centralizer and gives the margin formula on scan 127.
\end{proposition}
\begin{proof}
All maps are the quotients $q_U:G_{K'(X)}\to E_U$ induced by
the generic point. A section of $E_U$ lifts by (b); applying (a)
to the lift puts it in $D_v$ for a rational place of $X_{K'}$.
The map $q_U$ kills its inertia when the place belongs to $U$,
and identifies its decomposition group with the cusp group when
the place belongs to $X_{K'}\setminus U$. Thus its image is
respectively the point section or a member of that cusp packet.
This proves one implication.

For the reverse implication, a point section at $x\in U(K')$
lifts by any splitting of the $x$-adic decomposition sequence.
For a cusp $c$, the homomorphism $D_c(G_{K'(X)})\to D_c(E_U)$
is an isomorphism: it is an isomorphism on the inertia $T$ and
on the residue quotient $G_{K'}$. A given cusp section therefore
lifts through that actual isomorphism. Conjugators lift because
$q_U$ is surjective. This constructs every lift required by (b).

Now start with a birational section $s$. For every hyperbolic
open $U$, (F2) and Proposition~\ref{F-fixed-dictionary} give
its unique fixed pro-point in the completed universal cover of
$U$. For $U'\subset U$, forgetting coverings ramified in
$U\setminus U'$ maps that fixed pro-point to the unique one for
$U$. They are therefore compatible. Every finite separable
extension of $K'(X)$ is unramified over some such $U$, so these
compatible pro-points give points on the normalizations of $X$
in every finite extension of $K'(X)$. Their local rings form the
valuation ring of a prolongation of one place; equivalently this
is the usual inverse-limit description of places on a curve.
The section fixes all these points and hence preserves that
valuation. The place on $X$ is $K'$-rational, because its
decomposition image contains all of $G_{K'}$. This proves (a).
Two invariant prolongations would give two fixed pro-points at
some finite covering level, contradicting uniqueness. Thus the
comparison includes the uniqueness assertion that is valid in
one variable.
\end{proof}

\begin{proposition}[F.28 --- Proved: the proper-base part of the polycurve lifting reduction]
\label{F-proper-polycurve-lifting}
The universal curve-lifting assertion (F4)(b) over all finitely
generated extensions of $\mathbb Q$ implies (F4)(b) for every
hyperbolic polycurve whose tower is proper at every stage except
possibly the last. The comparison is given by taking the point
fibre of the projected section and then a valuation centered at
the resulting point or relative cusp. This proves the stated
part of the implication asked in F(e)(4); towers with a
cuspidal projected section on an affine intermediate base are
not included in the assertion.
\end{proposition}
\begin{proof}
Proposition~\ref{F-cuspidalization} proves the implication from
the universal curve-lifting assertion to the universal curve
point-or-cusp assertion. For a proper hyperbolic polycurve
$B/K$, induction shows that every section is a point section.
Indeed its projection to the preceding stage is, after geometric
conjugation, the section $s_b$ of a rational point $b$. Hoshi,
Proposition 2.4(ii), identifies the pullback
\[
 E_B\times_{E_{B'},s_b}G_K=E_{B_b}.
\]
This is a pullback along an actual point, with the same domain
$G_K$. The given section is now a section of the proper
hyperbolic fibre. The curve assertion makes it a point section,
completing the induction.

Apply this to the proper base $B$ of the last map $U\to B$.
The same pullback identifies the original section with a section
of $U_b$. It is either the section of a rational point or a cusp
section of this fibre. In the former case choose a regular system
of local parameters at the point of $U$. Their successive prime
ideals define a full flag; composing the resulting discrete
valuations gives a rank-$\dim U$ valuation of $K(U)$ with
residue $K$ and center at that point. A decomposition splitting
exists by Proposition~\ref{F-valuation-structure}; its image in
$E_U$ is the prescribed point section because finite \'etale
covers are unramified at that center.

In the cuspidal case take the smooth relative compactification
of $U\to B$. The relative cusp is a $K$-point $c$ over $b$,
and its boundary is \'etale over $B$. Let $z$ be a local
equation for that boundary and complete it by parameters of
$B$ at $b$. The same full-flag valuation has residue $K$ and
inertia $T^{\dim U}$. Its map to the local punctured tube is
the coordinate projection $T^{\dim U}\to T$ belonging to
$z$: on $n$th roots it is exactly the Kummer character of $z$.
Choose any splitting supplied by
Proposition~\ref{F-valuation-structure}. The prescribed fibre
cusp section differs from its image by a cocycle in
$Z^1(G_K,T)$. Insert that cocycle in the
$z$ coordinate of $T^{\dim U}$ and put zero in the others.
Twisting the decomposition splitting by it gives a section of
$G_{K(U)}$ whose image is the prescribed cusp section. This
constructs the required lift, rather than attempting to replace
a $G_K$-section by a section with domain $G_{K(B)}$.
\end{proof}

\begin{proposition}[F.29 --- Proved: point-theoretic maps and the proper-target reduction]
\label{F-point-theoretic}
Let $K$ be a number field or a finite extension of $\mathbb Q_p$,
$V/K$ a smooth geometrically connected variety, and $U/K$ a
hyperbolic curve. Call a continuous
$G_K$-homomorphism $f:E_V\to E_U$ \emph{point-theoretic}
when every point section over a finite extension becomes a point
section after a further finite extension. Such an $f$ is either
open or kills $P_V$ and is the point section of a unique
$K$-rational point composed with $E_V\to G_K$. In particular
every such $f$ comes from a unique morphism $V\to U$.

For proper $U$ over a number field $K$, condition (F1)(b) holds
for every continuous $G_K$-homomorphism $E_V\to E_U$. The all-source assertion
(F1) for this proper target and all finite extensions of $K$
is logically equivalent, via composition with point sections,
to (F2) for $U$ over those finite extensions. Thus the reduction
asked in F(b)(5) is proved for proper targets. For affine
targets the point-theoretic hypothesis is the additional
premise used in the proved map-level statement.
\end{proposition}
\begin{proof}
We use Carlson--Stix, Proposition C and Proposition 3.17 in its
hyperbolic-target range~\cite{F-CarlsonStix}: a point-theoretic
map between smooth curves with hyperbolic target is open or
contained in one point decomposition group. We checked its
argument and its field hypothesis: it is stated over every
Kummer-faithful field of characteristic zero. Both number fields
and finite extensions of $\mathbb Q_p$ have this property.
For the latter, a semiabelian variety over every finite extension
has a subgroup of finite index described by finite-rank valuation
groups and $p$-adic Lie groups; its intersection of all $n$-multiples
is zero. Equivalently one uses the compact open formal subgroup,
its $p$-power filtration, and the finitely generated valuation
quotient. Finite covers of a hyperbolic target have Euler
characteristic tending to minus infinity as their degrees grow.
The generalized-Jacobian dimension argument makes the induced
map on algebraic points have finite image. The group it generates
is both finitely generated and a quotient of a divisible
generalized-Jacobian group, hence zero. This proves constancy;
applying it at all finite covering levels gives the decomposition
group conclusion. The broader wording of the body proposition
only assumes nontrivial geometric fundamental group, but its
written Euler-characteristic argument requires hyperbolicity.
Our application has that hypothesis.

Suppose $f$ is not open. The induced map
$V(\overline K)\to U(\overline K)$ is well defined by point
section injectivity over finite extensions. Any two geometric
points of $V$ can be joined by a chain of curves: use affine
neighborhoods with nonempty successive intersections and
Bertini complete intersections through the chosen points,
then normalize. These curves and their points are defined over
some finite extension of $K$; puncturing further makes the source
curves hyperbolic. Composition with each curve's fundamental
group cannot be open, because that would make $f$ open. The
curve result therefore makes the point map constant on each
curve and then on all of $V(\overline K)$.

We now prove triviality on $P_V$, rather than infer it merely
from this constancy at the first level. After a finite Galois
extension $K'/K$, choose $x\in V(K')$, $y\in U(K')$ and
paths for which $f\circ s_x=s_y$ on $G_{K'}$. For each open
characteristic subgroup $N\subset P_U$, the subgroup
$H_N=N s_y(G_{K'})$ defines a geometrically connected finite
\'etale cover $U_N\to U_{K'}$. Its inverse image
$J_N=f^{-1}(H_N)$ defines a geometrically connected finite
\'etale cover $V_N\to V_{K'}$: it maps onto $G_{K'}$ since
it contains $s_x(G_{K'})$. The induced map $J_N\to H_N$
is still point-theoretic. Indeed a point section in the cover
maps to a point section downstairs, and its containment in
$H_N$ is exactly the fixed-coset criterion for a lift to
$U_N$, after a finite residue extension.

The same curve-chain argument makes
$V_N(\overline K)\to U_N(\overline K)$ constant. On the
fibres above $x$ and $y$ its map is the actual coset map
\[
 P_V/(P_V\cap J_N)\longrightarrow P_U/N,
 \qquad a\longmapsto f(a)N.
\]
This follows either from the finite covering fibre functors or
from their point-section stabilizers and injectivity on $U_N$.
Constancy therefore gives $f(P_V)\subset N$. Characteristic
open subgroups are cofinal because $P_U$ is topologically
finitely generated, so $f(P_V)=1$. Thus $f$ factors through a
section of $G_K$. That section restricts to $s_y$ over $K'$.
Normality of $G_{K'}$ and the self-normalizing property of
point decomposition groups make the same pro-point fixed by the
whole section; its image $y$
is $K$-rational, and the section is its point section. In the
open case Proposition~\ref{F-open-Hom} gives the unique dominant
morphism. The point case gives the unique constant morphism.

For proper $U$ over a number field, take an open subgroup
$J\subset E_V^{\circ}$
and a finite \'etale source cover with a closed point whose
cyclotomic section lands in $J$. Its composite under $f$ is
some section of the proper curve over a number field. The
proper case of Proposition~\ref{F-centralizer} makes its
geometric cyclotomic centralizer zero, whether or not that
section is known to come from a point. This kills the
centralizer of $f(J)$; projection to $G_K$ kills the arithmetic
one as in Proposition~\ref{F-geometric-necessary}. Thus (b)
is automatic. Finally (F2) for the proper target over every
finite extension makes the composite of $f$ with every point
section a point section. The first part then proves (F1).
Conversely the point-source instance of (F1) is (F2), as in
Proposition~\ref{F-point-equivalence}. These constructions
prove the stated equivalence of assertions without asserting
either unresolved universal surjectivity statement as a theorem.
\end{proof}

\begin{proposition}[F.30 --- Proved: the full curve-target reduction of (F1) to (F2)]
\label{F-full-curve-reduction}
Over all number fields, the following two universal assertions
are equivalent: (F2) for hyperbolic curves, and (F1) with a
hyperbolic curve as target and an arbitrary smooth connected
finite-type source. The comparison retains every non-open map
and every cusp packet. More precisely, the (F2) assertion implies
that a non-open arithmetic map to a hyperbolic curve has
geometric image either trivial or contained in one cusp inertia.
In the nontrivial case its entire image is in that cusp's
decomposition group, and condition (b) fails. This is a proved
implication between the assertions requested in F(b)(5), not
an assertion that their common universal surjectivity holds.
\end{proposition}
\begin{proof}
The reverse comparison is the point-source calculation of
Proposition~\ref{F-point-equivalence}. We prove the forward
comparison as a map between the two assertions. Work with a
continuous map $f:E_V\to E_U$ over a number field $K$, where
$V$ is smooth and geometrically connected. An open map is
handled by Proposition~\ref{F-open-Hom}, so consider a
non-open map. Choose a point $x\in V(K')$ after a finite
Galois extension $K'/K$. The (F2) assertion applied to
$s=f\circ s_x$ gives a fixed lifted interior point or cusp
$c$ in the completed universal cover of $U$. Write $I_c=1$
in the interior case and use the usual cusp inertia otherwise.

Choose a characteristic open subgroup $N_0\subset P_U$ whose
cover has compactification of genus at least two. Such a
geometric cover exists for every hyperbolic curve. In genus
zero use the degree-four Fermat covering of the tripod, after
selecting three punctures. In genus one an unramified cover of
the compactification first supplies at least two punctures,
and a double cover branched at two of them has genus two;
over an algebraically closed field the required square root of
their degree-two divisor exists in the Picard group. A
characteristic refinement preserves this genus bound. Crucially,
$N_0$ is chosen as an arbitrary geometric cover; it is not
required to contain $f(P_V)$.

For every characteristic open $N\subset N_0$, put
\[
 H_N=N s(G_{K'}),\qquad J_N=f^{-1}(H_N)
       \subset E_{V_{K'}}.
\]
These define geometrically connected finite \'etale covers
$U_N$ and $V_N$ over $K'$, since both subgroups project onto
$G_{K'}$. Their chosen point or cusp $c_N$ is fixed by $s$.
Let $X_N$ be the smooth compactification of $U_N$.
We claim that the homomorphism
\[
 J_N\longrightarrow H_N\longrightarrow E_{X_N}
\]
comes from the constant morphism $V_N\to X_N$ with value
$c_N$. The (F2) assertion for proper curves makes this map
point-theoretic, so Proposition~\ref{F-point-theoretic} gives
a unique morphism. To prove it constant, hold this $N$ and
this source $V_N$ fixed. Because $f(J_N)$ is non-open in
$H_N$, there are open subgroups $A\subset H_N$ containing
$f(J_N)$ with unbounded indices. Their covers $U_A$ remain
geometrically connected over $K'$. The compactifications obey
\[
 g(X_A)\geq1+[H_N:A]\bigl(g(X_N)-1\bigr),
\]
by Riemann--Hurwitz, hence have unbounded genus. Applying the
same proper point-theoretic argument gives morphisms
$h_A:V_N\to X_A$. A nonconstant such map is dominant and
induces an open fundamental-group map. Restriction in
degree-one $\ell$-adic cohomology is consequently injective:
a homomorphism from $P_{X_A}$ to $\mathbb Q_\ell$ vanishing
on an open subgroup vanishes everywhere. Therefore
\[
 2g(X_A)\leq \dim_{\mathbb Q_\ell}
       H^1(V_{N,\overline K},\mathbb Q_\ell),
\]
an impossibility for sufficiently large index. For such $A$
the map $h_A$ is constant. Its composition to $X_N$ induces
the original homomorphism, and uniqueness gives constancy of
the original morphism. Its value is $c_N$, by evaluating at
the distinguished lift of $x$.

Now evaluate that constant morphism on all geometric points
of $V_N$ above $x$. The map on the target compactified fibre
is, with the chosen universal paths, precisely
\[
 P_V/(P_V\cap J_N)\longrightarrow N\backslash P_U/I_c,
       \qquad a\longmapsto N f(a)I_c.
\]
One can check the formula on point stabilizers: the section
of the lift labelled by $a$ maps to
$f(a)s(G_{K''})f(a)^{-1}$ after a finite extension splitting
the fibres. Its point in $X_N$ is the lift labelled by
$Nf(a)I_c$. Constancy at every $N$ now gives
\[
 f(P_V)\subset\bigcap_{N\subset N_0}N I_c=I_c,
\]
because $I_c$ is closed and these characteristic subgroups
are cofinal. This argument includes both interior and boundary
fibres, with exactly the displayed quotient between them.

If $f(P_V)\ne1$, the inertia normalizer formula gives
$N_{E_U}(f(P_V))=D_c$. Since $f(P_V)$ is normal in $f(E_V)$,
the whole image is contained in $D_c$, which fails (b) by
Proposition~\ref{F-point-equivalence}. Otherwise $f$ factors
through a section of $G_K$. Condition (b) excludes its cusp
case, and the (F2) assertion makes it the section of a unique
point. This yields the desired constant morphism. A connected
regular source over $\mathbb Q$ reduces to the geometrically
connected case over its number field of constants; compatibility
over $G_{\mathbb Q}$ gives the required embedding of the
target constants. Thus the stated all-source comparison has
exactly the domain of (F1).
\end{proof}


\begin{proposition}[F.31 --- Proved using the supplied local theorem: the infinite-index dichotomy]
\label{F-local-dichotomy}
Let $K/\mathbb Q$ be finitely generated, $V/K$ a smooth
geometrically connected variety, $U/K$ a hyperbolic curve,
and $f:E_V\to E_U$ a
continuous $G_K$-homomorphism. Then exactly one of the following
three possibilities holds for its geometric image $H=f(P_V)$:
\[
 H\text{ is open in }P_U;\qquad H=1;\qquad
 1\ne H\subset I_c\text{ for one lifted cusp }c.
\]
In the third case $f(E_V)\subset D_c$ and the cusp is
$K$-rational. Consequently condition (F1)(b) implies that $H$
is open or trivial. In the trivial case the actual factorization
is
\[
 E_V\longrightarrow G_K\xrightarrow{\,s\,}E_U,
\]
with $s$ a continuous section. This proves the first assertion
of F(b)(5) with the local section theorem explicitly supplied
as tool (1) in the prompt~\cite{B_OpenAIPadic}. Its October
2026 preprint status
and lack of a formal verification are retained as the status
of that input; no independent proof of that input is asserted
here. Together with Proposition~\ref{F-full-curve-reduction},
this completes the requested curve-target reduction, without
asserting the remaining global section-surjectivity statement.
\end{proposition}
\begin{proof}
Choose an embedding $K\hookrightarrow k$, where $k$ is a
finite extension of some $\mathbb Q_p$. Such an embedding
exists: take a smooth model over the number field of constants,
a closed point after a finite extension, and an \'etale local
coordinate chart around it. Over a completion the chart contains
a $p$-adic open ball. Choose algebraically independent coordinates
in that ball; its lift is a generic point and gives an injective
evaluation of the function field into the completion. For a
number field one may simply take any finite-place completion.
Choose compatible separable closures. Base change gives the
actual cartesian
diagrams
\[
 E_{V_k}=E_V\times_{G_K}G_k,\qquad
 E_{U_k}=E_U\times_{G_K}G_k.
\]
The geometric groups are unchanged by the algebraically closed
extension in characteristic zero. No openness or injectivity
claim about $G_k\to G_K$ is needed. The induced $f_k$ therefore
has the very same geometric image $H\subset P_U$. Arithmetic
openness for a map over $G_k$ is equivalent to openness of its
geometric image.

The supplied local theorem, and its affine neighborhood
consequence (T1), say that every section of every hyperbolic
curve over every finite extension of $k$ is a point section or
a cusp section. For completeness, the non-open-image argument
in Proposition~\ref{F-full-curve-reduction} uses only this
local assertion and the point-theoretic statement of
Proposition~\ref{F-point-theoretic}, which was proved for
finite $p$-adic extensions as well. Here are its actual
constructions over $k$. After a finite extension choose
$x\in V(k')$ and let $c$ be the lifted point or cusp fixed
by $f_k s_x$. Choose an arbitrary characteristic target cover
whose compactification has genus at least two, and pull it
back to $V$; this fixes the source cover used next. Open
overgroups of the non-open image have arbitrarily large
indices. Their target compactifications have unbounded genus,
whereas a dominant map from that fixed source injects their
$H^1$ into its fixed finite-dimensional $H^1$. The local
proper section theorem makes all these maps point-theoretic,
so Proposition~\ref{F-point-theoretic} constructs them and
forces the maps of sufficiently high genus to be constant.
At every characteristic level $N$, the maps on the fibres
above $x$ are exactly
\[
 a\longmapsto N f_k(a) I_c\quad\text{in }N\backslash P_U/I_c.
\]
Their constancy proves $H\subset\bigcap_NNI_c=I_c$, where
$I_c=1$ for an interior point. This proves the local trichotomy
without needing condition (b) to remain true under localization.

Return to the global map. In the nontrivial inertia case,
$H\triangleleft f(E_V)$, and the geometric inertia-normalizer
identity is
\[
 N_{E_U}(H)=D_c.
\]
It follows that the entire global image lies in $D_c$. Its
projection is all of $G_K$, so the underlying cusp is
$K$-rational. Every element of $f(E_V^\circ)$ then centralizes
$I_c$, since the action on inertia is cyclotomic. Thus (b)
fails. In the remaining non-open case $H=1$, the quotient
property of $E_V\to G_K$ supplies the displayed continuous
map $s$, and the equation $p_Uf=p_V$ makes $s$ a section.
The cases are disjoint: a nontrivial cusp inertia has infinite
index in a hyperbolic geometric fundamental group. This last
fact also follows from its finite covers, whose Euler
characteristics cannot be those of a procyclic group.
\end{proof}

Statement B.39 extends this trichotomy to every sub-$p$-adic
field. Its explicit countable tower of finite gerbes descends
the homomorphism itself and then embeds that countable datum
in a finite $p$-adic field; the whole geometric homomorphism
and its peripheral subgroups are preserved. Statements
B.40--B.41 apply this extension to unramification over smooth
$p$-adic models and to the complete isotrivial section
classification. They retain the same supplied local input
and do not assert generic point realization for a nonconstant
relative section.

\begin{proposition}[F.32 --- Known, with exact hypotheses: higher-dimensional Hom and Isom ranges]
\label{F-higher-dimensional-ranges}
The cited higher-dimensional theorems yield the following precise
parts of F(b)(6).

For a sub-$p$-adic field $K$, a normal variety $V/K$, a hyperbolic
polycurve $U/K$ of dimension $n$, and an open arithmetic
homomorphism $f:E_V\to E_U$, Hoshi's Theorem A gives a unique
dominant morphism in each of these ranges:
\[
\begin{array}{c|l}
n&\text{additional hypotheses}\\ \hline
1&\text{none}\\
2&\ker f\text{ is topologically finitely generated}\\
3&\ker f\text{ is finite},\quad V\text{ is of LFG type},\quad\dim V\ge3\\
4&f\text{ is injective},\quad V\text{ is a hyperbolic polycurve},
       \quad\dim V\ge4 .
\end{array}
\]
Here LFG means that every nonconstant morphism from a normal
geometric variety to $V_{\overline K}$ has infinite geometric
fundamental-group image. In dimension three the resulting map
is quasi-finite and $\dim V=3$; in dimension four it is finite
\'etale and $\dim V=4$.

For $K/\mathbb Q$ finitely generated, Schmidt--Stix, Theorem 1.5
and Corollary 1.6, identify isomorphisms of strongly hyperbolic
Artin neighborhoods with outer arithmetic group isomorphisms.
The neighborhood hypothesis is a tower of elementary hyperbolic
fibrations in which every intermediate variety embeds locally
closed into a product of hyperbolic curves.

Holzschuh--Schmidt--Stix, in the sharp theorem at the end of
their paper, gives dominant relative Hom for $U\to S$ a
relative hyperbolic curve, with $S$ normal, connected,
excellent and of characteristic zero, and with sub-$p$-adic
generic field. The source $V$ is normal and connected,
$V\to S$ is flat and of finite type with connected geometric
generic fibre, and the smooth locus of $V\to S$ surjects onto
$S$. This includes the smooth surjective source in their
introductory Theorem A. Its input is a
$\pi_1$-open map of profinite \'etale homotopy types over
$\mathrm{Et}(S)$, with its commutation homotopy. When $S$ is a
hyperbolic polycurve this is exactly the corresponding relative
outer-homomorphism datum, through the canonical maps to
classifying spaces.

Every open arithmetic map between the objects of (F1) satisfies
(b). Thus the displayed Hom ranges prove both (a) and (b) for
their open maps. The Isom result is the isomorphism range of the
same comparison, through the explicit inclusion of isomorphisms
among open homomorphisms.
\end{proposition}
\begin{proof}
The hypotheses and conclusions are those of Hoshi, Theorem A,
Definition 2.5 and Corollaries 3.15--3.16\cite{F-HoshiPoly},
Schmidt--Stix, Definition 6.1 and Theorem 1.5/Corollary
1.6\cite{F-SchmidtStix}, and the attached August 2026
Holzschuh--Schmidt--Stix Theorem A and its sharpened final
theorem\cite{F-HolzschuhSchmidtStix}. These primary statements
were read directly, together with the final proof section.
For the last theorem, the smooth-locus surjectivity gives
the exact sequence
\[
 \pi_1(V_{\overline\eta})\longrightarrow E_V
     \longrightarrow E_S\longrightarrow1.
\]
The corresponding target sequence has injective geometric
fibre group, since that group is centre-free. For a relative
homomorphism the two image indices consequently agree:
\[
 [E_U:f(E_V)]
 = [\pi_1(U_{\overline\eta}):
                f_\eta(\pi_1(V_{\overline\eta}))].
\]
This is the explicit comparison between total and generic-fibre
openness. Their proof then constructs the generic-fibre map,
spreads the commutation homotopy by pro-\'etale descent, and
applies the extension criterion for hyperbolic curves. The
displayed flatness and smooth-locus hypotheses are retained.

For the asserted comparison of data, an \'etale $K(\pi,1)$
object has the canonical equivalence
$\mathrm{Et}(X)\simeq B E_X$. The universal property of
$1$-truncation identifies maps from any $\mathrm{Et}(V)$ into
$B E_U$ with homomorphisms of its fundamental groupoid. Doing
this in the slice over $B E_S$ retains the specified
commutation isomorphism; it is the relative groupoid datum,
with natural transformations as its conjugacies. Hyperbolic
polycurves are $K(\pi,1)$ by the successive homotopy fibre
sequences. Thus this supplies the precise morphism from the
relative homotopy theorem to the group formulation. When $S$
is not a $K(\pi,1)$, we retain the homotopy datum itself.

Finally let $J\subset E_V^\circ$ be open and $f$ open.
Then $f(J)$ is open in $E_U^\circ$. Its intersection with
$P_U$ is open in $P_U$ and has trivial centralizer there:
geometric hyperbolic polycurve groups are slim by the fibre
sequence and the surface-group statement. The arithmetic
projection of a centralizer element centralizes an open
cyclotomic kernel in $G_K$ and is trivial by the Hilbertian
centralizer theorem used in Proposition~\ref{F-geometric-necessary}.
The element therefore belongs to $P_U$ and is trivial. This
proves both readings of (b) for open maps without presupposing
the existence of a morphism. The cited reconstruction theorems
then give (a), and Proposition~\ref{F-geometric-necessary}
gives its necessary direction.
\end{proof}

\begin{proposition}[F.33 --- Proved and refuted: the first surface case and product targets]
\label{F-surface-products}
The assertion that an arbitrary open arithmetic map to a
hyperbolic polycurve surface must come from a dominant morphism
is false already for a smooth proper curve as source. For
$C/\mathbb Q$ a smooth proper curve of genus at least two,
there is a smooth geometrically connected ample curve
$D\subset C\times C$ for which
\[
 E_D\longrightarrow E_{C\times C}
\]
is surjective and satisfies (b), although its unique geometric
realization is the nondominant inclusion. Its kernel is not
topologically finitely generated.

For every product $U=U_1\times\cdots\times U_r$ of hyperbolic
curves, every open arithmetic map from a normal variety is
induced by a unique morphism to $U$, with dominance determined
by that morphism's actual image. On arbitrary arithmetic maps,
condition (b) is exactly the conjunction of the conditions (b)
on all coordinate maps. Consequently universal (F1) for product
targets over number fields is equivalent to universal (F2) for
curves over number fields, through those coordinate maps.
\end{proposition}
\begin{proof}
Choose a very ample line bundle on $C\times C$. Bertini gives
a nonempty open set of its sufficiently high-degree linear
system consisting of smooth geometrically connected divisors.
That open set has a rational point, since $\mathbb Q$ is
infinite; choose the corresponding $D$. For every connected
finite \'etale cover $Y\to(C\times C)_{\overline{\mathbb Q}}$,
the pullback of $D$ is ample and connected. Here connectedness
can also be checked directly by the Hodge index theorem:
a hypothetical disjoint decomposition into nonzero effective
divisors $D_1,D_2$ would give
$D_1D_2=0$ and $D_i^2=D_iD>0$, contradicting the single
positive direction of the intersection form. The finite-cover
criterion for fundamental groups therefore makes
$P_D\to P_{C\times C}$ surjective. Both arithmetic groups
surject onto $G_{\mathbb Q}$, so their arithmetic map is
surjective as well. It satisfies (b) by
Proposition~\ref{F-geometric-necessary}. The source has
dimension one and the target dimension two, so the inclusion
is not dominant: a dominant map would inject a function field
of transcendence degree two into one of degree one.
The product reconstruction below makes it the unique
realization of its group map. If the kernel were topologically
finitely generated, Hoshi's dimension-two theorem would instead
produce a dominant realization, a contradiction.

The exact arithmetic product formula is
\[
 E_U=E_{U_1}\times_{G_K}\cdots\times_{G_K}E_{U_r},
 \qquad P_U=\prod_iP_{U_i}.
\]
An open $f$ has open coordinate maps $f_i$. Mochizuki's
curve-target Hom theorem constructs their unique morphisms
$g_i$, and $g=(g_i)_i$ induces $f$ by this actual fibre-product
identity. Equality of all coordinates also proves uniqueness.
For every open $J\subset E_V^\circ$ one has
\[
 C_{P_U}(f(J))=\prod_i C_{P_{U_i}}(f_i(J)).
\]
The equivalence of geometric and arithmetic centralizer
readings proves the claimed coordinate test. Apply
Proposition~\ref{F-full-curve-reduction} to each coordinate
to obtain the implication from universal (F2) to universal
product-target (F1). Conversely take $U=X\times
(\mathbb P^1\setminus\{0,1,\infty\})$, multiply a noncuspidal
section of $X$ by the point section at $2$, and apply the
point-source instance of product-target (F1). Projection of
the resulting point gives the required point of $X$; cusp
packets are already supplied by the local construction.
This proves the reverse implication with the displayed maps.
\end{proof}

\begin{remark}[F.34 --- Open: the remaining surface Hom problem]
\label{F-surface-open}
After Proposition~\ref{F-surface-products}, the unrestricted
claim that openness implies dominance is refuted in the first
possible target dimension. The underlying reconstruction
question still left by F(b)(6) includes an open map to a
non-product hyperbolic polycurve surface whose kernel is not
topologically finitely generated, and for which the reconstructed
map to the base curve has no known flat presentation with
geometrically connected generic fibre whose smooth locus
surjects onto that base. This statement is
not declared proved or refuted here.

We tried the following explicit reduction. Compose with the
surface's base-curve projection and reconstruct that map by
Mochizuki. Over the largest open base where the source map is
smooth, use the fibre product
$E_{U\times_B B'}=E_U\times_{E_B}E_{B'}$ to form the
restricted relative datum. The missing input is openness on the
relative geometric fibre and extension across the removed base
points. Original openness can receive contributions from loops
around those removed points; it does not itself establish
openness of the image of the geometric generic-fibre group.
Hoshi's finite-generation condition supplies precisely a
normal-subgroup argument at this step. The HSS relative theorem
supplies an extension only under its displayed flatness,
smooth-locus surjectivity and homotopy hypotheses. No counterexample to
reconstruction in this remaining range has been constructed.
\end{remark}

\begin{proposition}[F.35 --- Proved comparisons: the 1981 and 1983 formulations]
\label{F-manuscript-comparison}
The comparisons requested in F(a) are given by the following
actual forgetful maps, restrictions and identifications.

For hyperbolic pairs $(X',S')$, $(X,S)$ over a number field,
put $U'=X'\setminus S'$ and $U=X\setminus S$. There is a
commutative square
\[
\begin{array}{ccc}
 \operatorname{Hom}^{\rm dom,\partial}(U',U)&\lhook\joinrel\longrightarrow&
                 \operatorname{Hom}^{\rm dom}(U',U)\\
 \downarrow\scriptstyle\sim&&\downarrow\scriptstyle\sim\\
 \operatorname{Hom}^{\rm open}_{\mathcal B}(\mathcal B_{U'},\mathcal B_U)
       &\lhook\joinrel\longrightarrow&
 \operatorname{Hom}^{\rm open}_{G_K}(E_{U'},E_U)/P_U .
\end{array}
\]
The boundary superscript means the literal boundary condition
$\overline f^{-1}(S)=S'$ of the corrected 1981 category.
The bottom arrow forgets the boundary diagram, its local
conjugators and its positive integer degrees. Its image is
exactly the inverse image, under the right-hand reconstruction,
of that boundary condition. Thus the open curve part of (F1)
implies the corrected dominant 1981 assertion; the latter
recovers the same (F1) assertion on precisely this displayed
image. Every dominant geometric map is represented in the
boundary category after the explicitly described shrinkings
in the proof. No equivalence obtained by silently forgetting
that boundary restriction is asserted.

On sections of a geometric curve, the corrected admissibility
condition of Longue Marche scan 79 is
\[
 \text{normalizes no cusp inertia}
 \quad\Longleftrightarrow\quad C_{P_U}(s(H_K))=1.
\]
Consequently its assertion that all admissible effective lifts
are point sections is equivalent to (F2), and also to the
point-source case of (F1), with every cusp packet retained on
the (F2) side. The letter's conjecture (7), with its packet
correction, is the same point-or-cusp assertion. Its
cyclotomic postscript is the right-hand zero-centralizer
condition in this displayed comparison.

The stronger guess of scan 433 holds for point sections and
fails for every cusp section: their $H_K$-fixed subgroups are
respectively $1$ and $I_c$. For the full Galois group and all
open subgroups its fixed subgroup is $1$ in both cases.
Thus the comparison from the full action to its restricted
$H_K$-action is literally restriction of the same action,
and the exact change of fixed subgroup is as stated.

Axioms (1)--(6) of scans 412--421, on realized geometric
objects with actual interior points, have the meanings and
descent maps proved in B.13. On those objects, replacing
``interior point'' by the intrinsic zero-centralizer selection
adds exactly the surjectivity assertion in (F2). For arbitrary
abstract outer actions, the inclusion of the geometric objects
into that larger collection does not prove essential
surjectivity or the intrinsic forms of axioms (5)--(6);
those are the unresolved tasks identified in Part B.

The first enlargement of the explicitly stated source category
in these 1983 assertions is the arbitrary regular finite-type
source in scan 48. The relation to the 1981 dimension-at-most-one
category is its restriction along
$(X,S)\mapsto X\setminus S$. A product of two positive-dimensional
hyperbolic curves has dimension two and is permitted as a source
in (F1), whereas it is not an object of that stated 1981 category.
All-morphism questions already occur in the 1981 text; their
mere presence in 1983 is not a new scope assertion.
\end{proposition}
\begin{proof}
The functor, its two-morphisms and the exact Kummer degree
condition are constructed in Propositions~\ref{A:typed-functor}
and~\ref{A:two-morphisms}. Positive peripheral degrees give
openness in the nonempty-boundary case by
Proposition~\ref{A:positive-degree-open}. The curve Hom theorem
and Proposition~\ref{A:relative-hom} give the two vertical
bijections. Proposition~\ref{A:boundary-recovery} proves that
the bottom forgetful arrow is injective and identifies its
image: a nontrivial peripheral image selects a unique target
inertia, its normalizer forces the local decomposition map,
and the Kummer formula recovers the positive integer degree.
The diagram therefore compares the actual functors of
Conjectures 5.5, its correction in section 9, and 7.1 with
(F1), without discarding their natural transformations.

Here is the asserted restriction for a dominant map whose
original punctures do not satisfy the boundary condition.
Extend it to the smooth complete curves and put
\[
 S_1=S\cup\overline f(S'),\qquad
 S'_1=\overline f^{-1}(S_1).
\]
The restrictions $X'\setminus S'_1\to X\setminus S_1$ and
the two open immersions give a commuting square of curves,
and hence a commuting square of their arithmetic groups.
The restricted map satisfies the exact boundary condition;
its exponents are its positive ramification indices. This is
the explicit relation between the unrestricted dominant map
and its boundary enhancement.

For the section comparison, every cusp-packet section
normalizes its inertia by construction, including all the
nongeometric tangential parameters. Proposition~\ref{F-centralizer}
proves the reverse recognition and the fixed-group formulas,
for sections and germs. Point-section and packet injectivity
make their images disjoint. Removing the whole union of cusp
packets from the equality asserted in (F2) gives exactly its
admissible-section assertion; restoring those packets gives
the converse. Proposition~\ref{F-point-equivalence} identifies
that same complement with (F1)(b) for a point source.
B.10 proves the full-action fixed-group statement. The
restriction to $H_K$ is the explicit homomorphism of actions
used in the display, so this also proves the requested
comparison with the stronger guess and the postscript.

For the axioms, B.13 uses finite base change for (1), deletion
and filling for (2), vanishing of fixed groups for (3), a
closed point after finite extension for (4), the equality
$(U\cup\{c\})(K)\setminus\{c\}=U(K)$ for (5), and
actual descent of the curve and its finite \'etale marked
subscheme for (6). Their uniqueness is supplied by the same
fixed-group vanishing on the configuration-space kernel.
These are maps on the realized subcategory. Extending them
to every abstract object requires precisely the realization
or surjectivity assertions still named in Part B, and is not
an additional theorem proved by this comparison. Finally the
source-category inclusion and the dimension of the displayed
product verify the claimed enlargement directly.
\end{proof}

\begin{proposition}[F.36 --- Known reductions and a proved obstruction diagram for F(c)]
\label{F-field-reduction}
Let $\mathrm{SC}_{\rm nf}$ denote (F2) over all number fields.
There is a proved implication from this assertion to (F2)
over finitely generated characteristic-zero fields for every
hyperbolic curve $U/K$ admitting a geometric finite \'etale
cover $Y\to U_{\overline K}$ and a nonconstant map
$Y\to A_{\overline K}$, where $A$ is defined over
$\overline{\mathbb Q}$. This includes all proper hyperelliptic
curves, every hyperbolic type of genus at most two, and every
affine curve with such a cover mapping nonconstantly to the
tripod. Every curve has a nonempty open subcurve in this last
range. This is an implication between the two section
assertions, not an assertion that $\mathrm{SC}_{\rm nf}$ has
been proved.

The general implication asked in F(c), with no cover
hypothesis, remains open in this return. The precise next
map after spreading already occurs over $K=k(B)$ for a
hyperbolic affine curve $B$ over a number field $k$:
For a smooth proper hyperbolic family
$X\to B$, a generic section gives
\[
 r:E_B\longrightarrow E_X,
 \qquad p r=\mathrm{id},
 \qquad E_{X_\eta}=E_X\times_{E_B}G_{k(B)}.
\]
Every such $r$ satisfies (F1)(b) and has non-open image. A
geometric reconstruction of this particular $r$ as a map
$B\to X$ is exactly a section of the family and supplies the
generic rational point. The curve-target dichotomy of
Proposition~\ref{F-local-dichotomy} does not provide that
higher-dimensional reconstruction.
\end{proposition}
\begin{proof}
The first implication and its ranges are Bresciani, Theorem B
and Corollaries A--C\cite{F-BrescianiFiniteType}. We checked
the proof: pass to a finite cover and field extension carrying
the map to a constant abelian variety; spread over a base
curve; use the coherence of the constant abelian family and
pull it back along a finite morphism by Hilbert irreducibility.
A section neighborhood whose compactification has no rational
point supplies the contradiction. The cover-to-tripod range
maps after a further cover to the constant elliptic curve of
$j=0$. The general Theorem A in that paper also uses weak
Bombieri--Lang; that additional assertion is not taken as an
established input here. Sa\"idi's earlier relative-Jacobian
approach has an additional Tate--Shafarevich-module condition;
we do not remove it by renaming the specialization problem.

For the displayed map, Bresciani's Corollary 3.4 was checked
with its proof. Inertia at a point of $B$ would give a
cyclotomic map $T\to P_{X_b}$. Its image can be detected in
the abelianization of a suitable finite cover, where weights
$-2$ and $-1$ exclude it. Thus the generic section descends
to the actual splitting $r$ of the relative fundamental-group
sequence. The fibre-product formula is the smooth hyperbolic
fibre sequence, so base change back to $G_{k(B)}$ recovers the
original section.

To prove (b), let $J\subset E_B^\circ$ be open. A geometric
centralizer of $r(J)$ projects to a centralizer of $J$ in
$P_B$, which is zero by the open case. It consequently lies
in the fibre group $P_{X_{\overline\eta}}$. Pull $J$ back
under $q:G_{k(B)}\to E_B$. Its intersection with the
cyclotomic kernel contains $H_L$ for a finite extension
$L/k(B)$. The proper-curve case of
Proposition~\ref{F-centralizer}, applied to the generic
section over $L$, kills the remaining centralizer. The
arithmetic centralizer is killed by the same projection
argument and the arithmetic centralizer theorem. On the
other hand $r(E_B)$ meets the infinite fibre group trivially,
so it has infinite index in $E_X$.

A morphism $g:B\to X$ inducing $r$ has $p g$ inducing the
identity on $E_B$. Curve Hom injectivity gives $p g=\mathrm{id}_B$,
and its generic point is the required $k(B)$-point. In the
other direction a generic rational point extends to a section
of this proper family over the regular curve $B$ by the
valuative criterion; its group map $r_g$ satisfies $r_gq=rq$ on $G_{k(B)}$,
and $q$ is surjective, so $r_g=r$. Thus this identifies
the missing reconstruction precisely.

Localizing the generic field in a finite $p$-adic extension
turns its section into a point by the supplied local theorem.
It gives a point of $X(k')$, not a proof that the point lies
in $X(K)$. The field inclusion induces the actual injection
$X(K)\to X(k')$; even on the tripod this map is not onto,
as a $p$-adic element transcendental over the countable field
$K$ gives a missing point. No extension of its section back
to $G_K$ is claimed in this example. It shows exactly why
this base-change step alone has not constructed the desired
rational point. The unresolved compatibility is the family
section reconstruction just displayed.
\end{proof}

\begin{proposition}[F.37 --- Known and proved comparison: the birational part of field reduction]
\label{F-birational-field-reduction}
The number-field assertion (F2) implies the one-variable
birational assertion (F4)(a) over every finitely generated
field of characteristic zero. Therefore, under this exact
implication of assertions, every section over such a field
which already has a birational lift is a point section or
cuspidal. The remaining lifting map is still the quotient
$G_{K(U)}\twoheadrightarrow E_U$; the comparison does not
manufacture a lift of an arbitrary section.
\end{proposition}
\begin{proof}
Proposition~\ref{F-birational-equivalence} sends the universal
number-field (F2) assertion to its one-variable birational
assertion by the compatible fixed-pro-point construction over
all open subcurves. Sa\"idi--Tyler, Theorem C
\cite{F-SaidiTyler}, sends the birational assertion over number
fields to the same assertion over all finitely generated
extensions of $\mathbb Q$. We opened their primary paper and
checked the exact quantifier and the reduction in section 5.3:
reduce to the projective line; select a characteristic cover of
a constant punctured projective line whose compactification is
hyperbolic; form the subgroup generated by that characteristic
geometric subgroup and the given section. This supplies an
isotrivial hyperbolic birational neighborhood. The isotrivial
relative-Jacobian theorem supplies the needed finite
Tate--Shafarevich group in their argument. Their Theorem C
therefore has no residual Tate--Shafarevich assumption.
Finally a birational lift belongs to the rational-place
decomposition group by that assertion, and its image in
$E_U$ is the point or cusp section according to the location
of the place. These are precisely the quotient maps used in
Proposition~\ref{F-birational-equivalence}.
\end{proof}

\begin{proposition}[F.38 --- Known, with checked proof: what inductive cuspidality gives]
\label{F-inductive-cuspidality}
Let $X=X_n\to\cdots\to X_0=\operatorname{Spec}K$ be a
hyperbolic polycurve over a finitely generated characteristic-zero
field. Write $\mathcal S_X$ for its section classes and
$C_{X,f}$ for the recursively defined set containing every
section whose projection is cuspidal, together with fibre
cuspidal sections above rational base points. The universal
curve section assertion over $K$ is equivalent to independence
of $C_{X,f}$ from the chosen tower, with the known
point/packet injections understood. For a normal-crossings
compactification $X\subset Z$, let $C_{X,Z}$ be the sections
coming from punctured henselian neighborhoods of rational
boundary points, and let $C_X$ be their closure. The checked
comparison is
\[
 C_X\subset\overline{C_{X,f}}.
\]
It gives no reverse inclusion or birational lift of a section
merely because its projection to an affine base is cuspidal.
Thus it does not complete the remaining case of F(e)(4).
\end{proposition}
\begin{proof}
These are Carlson--Corwin's Proposition B and the proposition
labelled \texttt{prop:fibinclusion} in the attached November
2026 source\cite{F-CarlsonCorwin}; we read the entire
fibration section and its proofs. The forward implication is
induction using the actual point-fibre pullback: a section
whose base projection is not in the recursively selected cusp
set belongs to a fibre over a rational point. The curve
assertion then exhausts that fibre's sections. Consequently
$C_{X,f}=\mathcal S_X\setminus X(K)$, independently of $f$.
For the converse take a hyperbolic curve $B$ and the tripod
$T$, and use the two projections of $B\times T$. Their
section-set product and the exact two recursive subsets are
\[
\begin{split}
 \mathcal S_{B\times T}&=\mathcal S_B\times\mathcal S_T,\\
 C_{B\times T,\operatorname{pr}_B}
   &=(C_B\times\mathcal S_T)\cup(B(K)\times C_T),\\
 C_{B\times T,\operatorname{pr}_T}
   &=(\mathcal S_B\times C_T)\cup(C_B\times T(K)).
\end{split}
\]
For $s\in\mathcal S_B$ and any $c\in C_T$, equality of the
last two subsets puts $s$ in $B(K)\cup C_B$, proving the
converse with its actual projection maps.

For the inclusion, choose a compatible smooth compactification
of the tower by resolving the proper compactification of the
last fibration. A rational boundary point maps either to a
boundary point of the compactified base, or to a rational
interior point of the base. In the first case its section
projects into the recursively selected boundary set. In the
second it lies on a relative cusp, and the punctured local
fundamental group is that of the cusp in the fibre: the
relative boundary is \'etale with multiplicity one. This
gives $C_{X,Z}\subset C_{X,f}$ for that compactification;
taking closures gives the display. The proof begins with an
actual local boundary section and maps it into the recursive
set. Starting with an arbitrary section in the first term
$f_*^{-1}(C_B)$ supplies none of that local data. The
unconstructed reverse step is exactly the one needed to
extend Proposition~\ref{F-proper-polycurve-lifting} to an
affine intermediate base.
\end{proof}

\begin{proposition}[F.39 --- Proved: explicit parameter transport and its GT-coordinate formula]
\label{F-tangent-transport}
Fix the rational tangent of parameter $1$ at $0$ on
$T=\mathbb P^1\setminus\{0,1,\infty\}$. Every parameter
$t\in\mathbb Q^*$ has the following exact characterization
using a specified four-puncture diagram. Put
\[
 W_t=\mathbb P^1\setminus\{0,1,t,\infty\},\qquad
 p(z)=z,\quad q(z)=z/t,
 \qquad p,q:W_t\longrightarrow T.
\]
Repeated punctures at $t=1$ are simply omitted. There is a
unique class in the cusp packet at $0$ of $W_t$ whose image
under $q_*$ is the parameter-$1$ class; its image under
$p_*$ is the parameter-$t$ class of $T$. In packet
coordinates the actual maps are
\[
 p_*:c\longmapsto c,\qquad
 q_*:c\longmapsto c-\kappa(t).
\]

In a normalized tripod presentation
$P_T=\langle x,y\rangle^{\wedge}$ with the parameter-$1$
section and
\[
 \rho_\sigma(x)=x^{\chi(\sigma)},\qquad
 \rho_\sigma(y)=f_\sigma^{-1}y^{\chi(\sigma)}f_\sigma,
 \qquad f_\sigma\in[P_T,P_T],
\]
changing that section by a cusp cocycle $c$ gives exactly
\[
 f^{\,c}_\sigma=f_\sigma x^{-c(\sigma)},\qquad
 (f^{\,c}_\sigma)^{\rm ab}=(-c(\sigma),0).
\]
Thus inserting the already specified Kummer cocycle
$\kappa(t)$ into these relations recognizes its parameter
class. Neither construction removes its input: the first
retains the four-puncture diagram with cross-ratio $t$, and
the second retains the Kummer character of $t$ and the
chosen origin. They do not yet answer the intrinsic
single-tripod GT-recognition question left in
Proposition~\ref{F-open-tangents}.
\end{proposition}
\begin{proof}
At $0$ both maps have ramification degree one. The first local
coordinate has leading coefficient $1$, and the second has
leading coefficient $1/t$. The tangent convention of
Proposition~\ref{F-torsors} therefore gives the two displayed
Kummer maps. In particular $q_*$ identifies the actual local
decomposition extensions, on both their inertia and their
$G_{\mathbb Q}$ quotient. Its inverse carries the chosen
parameter-$1$ section to the unique local section with class
$c=\kappa(t)$. The cusp-packet injection then gives uniqueness
in the ambient section set as well.

For the second formula write the new section as
$s_c(\sigma)=x^{c(\sigma)}s_1(\sigma)$. Its action is
$\operatorname{Int}(x^{c(\sigma)})\rho_\sigma$. It has the
same image of $x$, and its image of $y$ is
\[
 x^{c(\sigma)}f_\sigma^{-1}y^{\chi(\sigma)}
                   f_\sigma x^{-c(\sigma)}
 = (f_\sigma x^{-c(\sigma)})^{-1}y^{\chi(\sigma)}
                   (f_\sigma x^{-c(\sigma)}).
\]
This proves the exact transformation. Taking abelianization
proves the last formula. Existence of an inertia conjugation
for which $f^{\,c}_\sigma x^{\kappa(t)(\sigma)}$ has zero
abelianization for every $\sigma$ is precisely
$[c]=[\kappa(t)]$; changing roots changes the cocycle by the
same coboundary. This is a useful coordinate test, but its
explicit input character is exactly the data that an intrinsic
recognition rule would have to recover, so no stronger status
is assigned to it.
\end{proof}

\begin{proposition}[F.40 --- Refuted, with checked construction: higher-rank valuative uniqueness]
\label{F-valuation-nonuniqueness}
The uniqueness of a valuation normalized through its inertia,
and the inference from inertia normalization to preservation
of that valuation, fail even with nonzero inertia and residue
field exactly $\mathbb Q$. There are distinct valuations
$v_0,v_1$ of $L=\mathbb Q(z,x,y)$ and one section
$s:G_{\mathbb Q}\to G_L$ such that
\[
 \Gamma_{v_a}=\mathbb Z\oplus_{\rm lex}\mathbb Q,
 \quad\kappa(v_a)=\mathbb Q,
 \quad I_{v_0}=I_{v_1}=T\ne1,
 \quad s(G_{\mathbb Q})\subset N_{G_L}(I_{v_0})
                             =N_{G_L}(I_{v_1}),
\]
but $s(G_{\mathbb Q})\not\subset D_{v_1}$. Accordingly the
higher-dimensional version of (F4)(a) must retain actual
valuation preservation. Its written uniqueness with respect to
inertia normalization is false; the example does not assert
nonuniqueness among valuations actually preserved by this section.
\end{proposition}
\begin{proof}
This is F149-ZD5, whose full construction F149-D3 we checked.
Let $w$ be the $z$-adic valuation, with residue field
$\mathbb Q(x,y)$. In the Hahn field $\mathbb Q((t^{\mathbb Q}))$
define two restrictions $r_a$, for $a=0,1$, by
\[
 x-a=t,\qquad y=2+\sum_{j\ge1}t^{j+1/j!}.
\]
Their residue field is $\mathbb Q$ and their value group is
$\mathbb Q$. To verify the latter, form the norm polynomial
of $Y-2-\sum_{j\le m}X^{j+1/j!}$ from
$\mathbb Q(X^{1/m!})$ to $\mathbb Q(X)$. At the displayed
series its identity factor has value
$m+1+1/(m+1)!$, and the other factors have total value in
$(1/m!)\mathbb Z$. Induction therefore puts $1/(m+1)!$ in
the value group. Conversely every series exponent of a
rational expression belongs to the union of those groups.
Unbounded value-group index proves that $y$ is transcendental
over $\mathbb Q(t)$, so this really defines valuations of
$\mathbb Q(x,y)$. Their centers are $(0,2)$ and $(1,2)$.

Compose each $r_a$ with $w$. The divisible final convex
subgroup is $\mathbb Q$, so
Proposition~\ref{F-valuation-structure} gives
$I_{v_a}=I_w=T$ and $N(I_{v_a})=D_w$, as actual subgroups
for prolongations refining the same $w$. Construct a section
of $G_{\mathbb Q(x,y)}\to G_{\mathbb Q}$ along the rational
flag $y=2,x=0$, then lift it through a splitting of the
$w$-decomposition sequence. Its image lies in $D_w$ and
normalizes both inertia groups.

Restricting this section to $\mathbb Q(x)$, then to the
fundamental group of
$\mathbb P^1\setminus\{2,3,\infty\}$, gives the point
section at $0$. Containment in $D_{v_1}$ would instead give
the point section at $1$, since $r_1(x-1)>0$. Their distinctness
follows from the injective point-section map; equivalently
use its two Kummer coordinates. Thus the containment fails.
This verifies both refutations by the stated restriction maps.
\end{proof}

\begin{proposition}[F.41 --- Refuted descent without continuity: an explicit abstract action]
\label{F-discontinuous-action}
There is an abstract semilinear action of $G_{\mathbb Q}$ on
an algebraic universal covering $\mathcal U$ of the tripod
which is not functionwise continuous on any open subgroup.
It admits no finite-type quotient $\mathcal U\to W$ which
is a connected pro-\'etale covering and whose induced
semilinear action is functionwise continuous on an open
subgroup of $G_{\mathbb Q}$. Thus the finite-form descent
step in Proposition~\ref{F-continuous-action} is false
without continuity.

This action nevertheless has the unique fixed pro-point
asserted in (B$'$), even for $H_{\mathbb Q}$. It proves
that the forgetful map from the continuous actions to
abstract actions misses actions on this same covering, and
that the specified finite-form reduction does not extend
to them. Fixed-point existence for every abstract action
is not proved or refuted by this example.
\end{proposition}
\begin{proof}
Choose distinct odd primes $p_j$. Restriction to the
multiquadratic extension gives a continuous surjection
\[
 q:G_{\mathbb Q}\longrightarrow
                  A=\prod_{j\ge1}\mathbb F_2.
\]
The independence follows from unique prime factorization:
no nonempty finite product of these primes is a rational
square. Let $D=\bigoplus_{j\ge1}\mathbb F_2\subset A$.
Choose an $\mathbb F_2$-linear functional
$\lambda_0:A/D\to\mathbb F_2$ which takes the class of
$(1,1,\ldots)$ to $1$, by extending that nonzero vector
to a basis. Put $\lambda=\lambda_0q$, with the quotient
by $D$ understood. Its kernel is dense because it contains
the inverse image of the dense subgroup $D$ and $q$ is
an open quotient. It is not open: a continuous functional
on $A$ vanishing on $D$ vanishes identically. Its
restriction to every open subgroup of $G_{\mathbb Q}$
is nonzero and has dense kernel.

Put $T=\mathbb P^1\setminus\{0,1,\infty\}$ and fix the
rational interior point $x=1/2$. Choose a geometric point
above it on the universal covering and write $s$ for the
resulting continuous point action of $G_{\mathbb Q}$.
The rational involution $a:z\mapsto1-z$ has a unique lift
$\widetilde a$ fixing the chosen pro-point: two such lifts
differ by a deck transformation fixing a geometric fibre
point, hence agree. Its square is the identity, and it
commutes with $s(G_{\mathbb Q})$, since conjugating the
lift by $s(\sigma)$ gives another lift of the same rational
involution fixing the same point. Define
\[
 r(\sigma)=\widetilde a^{\lambda(\sigma)}s(\sigma).
\]
Commutation proves directly that $r$ is an action. It is
semilinear with the original action on $\overline{\mathbb Q}$.
On the rational coordinate $z$ it acts by $z$ or $1-z$,
so its stabilizer is $\ker\lambda$; this is not open even
after restriction to any open subgroup. Every quadratic
extension of $\mathbb Q$ lies in a cyclotomic extension
(for these prime radicals one can use the quadratic Gauss
sum in $\mathbb Q(\mu_{4p_j})$). Thus $\lambda$ vanishes
on $H_{\mathbb Q}$ and $r|_{H_{\mathbb Q}}=s|_{H_{\mathbb Q}}$.
Proposition~\ref{F-fixed-dictionary} gives its unique fixed
pro-point, which is also fixed by $r(G_{\mathbb Q})$.

We prove the stronger quotient assertion, rather than
just exhibit a noncontinuous coordinate. Write $\Omega$
for the function field of $\mathcal U$. Suppose that an
open subgroup $G'\subset G_{\mathbb Q}$ preserves a
finite-type quotient field $M\subset\Omega$ and that
its induced action on $M$ is functionwise continuous.
For $f\in M$, the two maps
\[
 G'\longrightarrow\Omega,\qquad
 \sigma\longmapsto r(\sigma)f,\quad
 \sigma\longmapsto s(\sigma)f
\]
are continuous when $\Omega$ is discrete. They agree on
the dense subgroup $G'\cap\ker\lambda$, so they agree
everywhere. Choose $\sigma\in G'$ with $\lambda(\sigma)=1$.
Commutation then gives $\widetilde a f=f$. Consequently
$M\subset\Omega^{\widetilde a}$.

If $\mathcal U\to W$ were a connected pro-\'etale
covering with function field $M$, $\widetilde a$ would
be a deck automorphism fixing its chosen geometric point.
The action of the deck group of a connected pro-\'etale
cover on a geometric fibre is free: this is the usual
unique path-lifting statement on every finite \'etale
level, followed by the inverse limit. Hence
$\widetilde a=1$, contradicting its action $z\mapsto1-z$.
No such quotient exists. This gives the exact obstruction
to the specified reduction and leaves the universal
abstract fixed-point assertion with its stated open status.
\end{proof}

\begin{proposition}[F.42 --- Proved: lifting above a base cusp with good marked reduction]
\label{F-good-boundary-lifting}
Let $U\to B$ be a hyperbolic curve fibration over a hyperbolic
curve over a finitely generated characteristic-zero field $K$.
Let $c\in(\overline B\setminus B)(K)$. Suppose that near $c$
the fibration has a smooth proper relative compactification
with finite \'etale boundary. Write $Y_c$ for its hyperbolic
special fibre with that boundary deleted.

There is the following exact extension comparison, after a
choice of specialization paths:
\[
 E_U\mathbin{\times}_{E_B}D_c
       \simeq E_{Y_c}\mathbin{\times}_{G_K}D_c.
 \tag{F.42.1}
\]
Consequently every section of $E_U\to G_K$ whose base
projection is in the packet at $c$ corresponds to a pair
consisting of that same cusp section and a section of
$E_{Y_c}\to G_K$. When the latter is a point or cusp
section, the original section has a birational lift to
$G_{K(U)}$, constructed by a rank-two valuation and its
specified inertia coordinates.

In particular the universal curve-lifting assertion (F4)(b)
over finitely generated fields implies (F4)(b) for all such
surfaces whose fibration has good marked reduction at every
rational base cusp after every finite extension. This extends
Proposition~\ref{F-proper-polycurve-lifting} to this affine-base
range. The remaining affine-base case includes genuinely bad
marked reduction at the cusp selected by the projected section.
\end{proposition}
\begin{proof}
Choose a parameter $t$ at $c$ and a prolongation of its
valuation. The cusp decomposition extension $D_c$ identifies
with $G_{K((t))}$: both have kernel $T$ and quotient $G_K$,
and compatible roots of $t$ identify their kernels and the
uniformizer splittings. The comparison of the global and
local relative fundamental-group sequences gives
\[
 E_{U_{K((t))}}\xrightarrow{\sim}
            E_U\mathbin{\times}_{E_B}G_{K((t))}.
\]
On the fibre group this map is the characteristic-zero
geometric base-change isomorphism; on the quotient it is the
identity. The relative fibre sequence is injective for a
hyperbolic fibration, so these two facts prove the displayed
isomorphism of extensions.

Marked good reduction identifies that fibre group with
$P_{Y_c}$ and makes its outer $G_{K((t))}$-action the
composite of the specialization map to $G_K$ with the
outer action of $Y_c$. This is the smooth marked
specialization comparison checked in
Proposition~\ref{F-reduction}; equivalently it follows from
the Kummer-\'etale comparison for the smooth pair, with
its horizontal boundary \'etale over the trait.
Here the residue characteristic is zero, so all geometric
finite covers are in its tame range.

We spell out why outer action suffices at this point.
For a centre-free profinite group $P$, an extension $E$
of a profinite group $G$ by $P$ is identified with
\[
 E\longrightarrow
 \operatorname{Aut}(P)\mathbin{\times}_{\operatorname{Out}(P)}G,
 \qquad e\longmapsto(\operatorname{Ad}(e)|_P,p(e)).
\]
Its kernel is $Z(P)=1$. For surjectivity, choose a lift of
the specified element of $G$ and correct its automorphism
by the unique required element of $P$. Compactness gives a
homeomorphism. Applying this to the centre-free group
$P_{Y_c}$ and the just identified outer action proves
(F.42.1), as an actual map of extensions after the stated
path choices. A section whose projection is the chosen
cuspidal section lands in the left side, and its two right
components give exactly the asserted pair. Conjugacies by
the geometric fibre act simultaneously through this map.

Suppose first that the residue section is a point section
at $y\in Y_c(K)$. In the smooth surface model around $y$,
complete $t$ by a local fibre parameter $z$. The flag
valuation with successive parameters $t,z$ has value group
$\mathbb Z\oplus_{\rm lex}\mathbb Z$ and residue field
$K$. Its decomposition sequence splits by
Proposition~\ref{F-valuation-structure}. In its image in
$E_U$, the $z$-inertia is killed, since $y$ is an interior
point of the special fibre, while the $t$-inertia is exactly
the base cusp inertia. Adjust a chosen splitting by the
cocycle in its $t$-coordinate that gives the prescribed
base packet member. The image of this adjusted splitting
has the specified two components in (F.42.1), hence is the
original section.

If the residue section is cuspidal at a relative cusp
$y\in(\overline Y_c\setminus Y_c)(K)$, choose $z$ as the
local equation of that horizontal boundary. The two
boundary divisors have normal crossings, and the local
Kummer maps send the inertia of the same rank-two valuation
to the $t$ and $z$ inertia coordinates. Each is a copy of
$T$, with cyclotomic action. The prescribed two packet
members differ from the images of a chosen decomposition
splitting by two cocycles in $Z^1(G_K,T)$. Twisting that
splitting by this ordered pair gives their simultaneous
lift. These maps are the same coordinate root maps used
in Proposition~\ref{F-proper-polycurve-lifting}; no section
with domain $G_{K(B)}$ has been inferred from one with
domain $G_K$.

Finally universal curve lifting gives the universal curve
point-or-cusp assertion by
Proposition~\ref{F-cuspidalization}. Project a section to
$B$. A point projection is handled by its actual fibre
and the full-flag construction, without a properness
requirement on $B$. A cusp projection is handled by
(F.42.1) and the preceding construction. This proves the
stated implication of lifting assertions in the additional
range. At bad marked reduction the residue extension of
Proposition~\ref{F-reduction} need not be that of a smooth
hyperbolic curve; indeed its checked equivalence excludes
such a realization after every finite extension when
potentially good marked reduction fails. Thus applying
curve (F2) to an imagined smooth residue fibre would not
supply the remaining step. A fixed component or compatible
logarithmic valuation for that extension is still required.
\end{proof}

\subsection{The added LM5 comparison: tubes, inertia, and actual sections}

\begin{proposition}[F.43 --- Proved: the normal-tube comparison with FLAGGED\_149]
\label{F-normal-tube-comparison}
Let $(X,D)$ be a smooth normal-crossings pair over a characteristic-zero
field, $U=X\setminus D$, and $S=D_J^*$ a connected open stratum with
its $r$ branches labelled. Write $N_j$ for their normal lines on $S$,
$P_S=\prod_{j\in J,S}N_j^\times$, and $\mathcal T_S$ for the
punctured henselian normal tube. Paths and geometric points are fixed
in every displayed fundamental-group map.
\begin{enumerate}
\item Normal specialization identifies the finite-covering categories
of $\mathcal T_S$ and $P_S$. Its profinite geometric homotopy sequence is
\[
 \pi_2^{\rm et}(S_{\bar K})\xrightarrow{(c_1(N_j))_j}
 T^J\longrightarrow\pi_1(\mathcal T_{S,\bar K})
 \longrightarrow\pi_1(S_{\bar K})\longrightarrow1.
 \tag{F.43.1}
\]
The first map means Chern transgression, evaluated on spherical
classes. The composite $\iota:T^J\to P_U$ kills its image. Thus
injectivity of the full global inertia implies left exactness in
(F.43.1). For a relative cusp of a hyperbolic smooth pair, the tube
extension maps isomorphically onto the cusp normalizer extension
of Proposition~\ref{F-torsors}, with extension class $c_1(N)$.

\item At a geometric boundary point, $\iota:T^J\to P_U$ is injective
precisely when every finite covering germ is an open and closed
sub-germ of the restriction of a finite global covering. If $R_{\bf n}$
is the finite root stack with orders $n_i$, and $M_{\bf n}$ is the
closed normal subgroup of $P_U$ generated by the $n_i$th powers of
boundary inertia, then
\[
 \pi_1(R_{{\bf n},\bar K})=P_U/M_{\bf n},\qquad
 \ker(T^J\to\pi_1(R_{{\bf n},\bar K}))
       =\iota^{-1}(M_{\bf n}),\qquad
 \bigcap_{\bf n}\iota^{-1}(M_{\bf n})=\ker\iota.
 \tag{F.43.2}
\]
In particular a finite root stack can enlarge this kernel.

\item Horizontal cusp inertia injects in a tower of punctured
hyperbolic smooth pairs. In the recursively compatible stable
toroidal range of LM5 F.39, the full crossing inertia injects as
well. The hypothesis here retains the stable incidence data:
the kernel of the local contact matrix is the fibre puncture
or the essential vanishing cycle. On split branches the Galois
action is cyclotomic; in general it also permutes the factors.
A monomial map $z'_j=u_j\prod_i z_i^{m_{ji}}$ induces the matrix
$(m_{ji})$ on these inertia groups.

\item For $y^2=x^3+t$, the model
\[
 [\{Y^2=X^3+1\}\times\mathbb A^1_s/\mu_6],\qquad
 \zeta(X,Y,s)=(\zeta^{-2}X,\zeta^{-3}Y,\zeta s),
 \tag{F.43.3}
\]
has open part identified with the original family by
$t=s^6$, $x=s^2X$, $y=s^3Y$. Its normal meridian injects;
the full corner meridian does too when the infinity section
is retained. On the ordinary embedded cusp resolution the
two exceptional-tail crossing maps instead have respective
kernels $\widehat{\mathbb Z}(-3,1)$ and
$\widehat{\mathbb Z}(-2,1)$. These are maps into the same
open fundamental group, with the explicit substitutions above
and in the proof.

\item The maps in (F.43.1), restricted to a relative cusp or a
trait meeting that cusp, are exactly the maps in FLAGGED\_149
C3, C4, C6, C8, and C9. C1, C2, and C5 give the corresponding
extension tests for an \emph{already given rational geometric
section}. Their geometric premise is retained; the tests do not
construct a geometric point from an arbitrary section.
\end{enumerate}
This supplies the added comparison of LM5 F.21--F.23, F.35,
F.39 with FLAGGED\_149 C1--C6, C8, C9.
\end{proposition}
\begin{proof}
We read the cited return proofs directly. Here are the actual
comparison maps and the kernel checks, so that equality of names
for an inertia group is not used in place of a comparison.

In a strict normal-crossings chart the covers are dominated by
adjoining $z_j^{1/n}$. On an overlap
$z'_j=u_jz_{\sigma(j)}$, their root coordinates are identified by
roots of $u_j|_S$. These are the transition torsors of the normal
lines. Normalization after the root adjunction is unramified,
and descent for finite covers makes this local construction an
equivalence with the covering category of $P_S$. The same
transition data in the inverse system of root gerbes gives the
profinite homotopy comparison: its fibre is $BT^J$, its twisting
classes are the Kummer classes of $N_j$, and its boundary is
their Chern transgression. This is the normal-bundle instance
of the log comparison of Carchedi--Scherotzke--Sibilla--Talpo,
Theorem A and Corollary 5.5 \cite{F-CSST}; their fine saturated,
log-regular, locally noetherian hypotheses hold for this smooth
normal-crossings pair. In characteristic zero the comparison
uses all primes. This justifies the higher homotopy sequence,
including its finite-covering truncation.

Restriction of a global cover to the tube is the functor
giving $\pi_1(\mathcal T_S)\to\pi_1(U)$. Its composite with
the fibre map kills the Chern boundary by (F.43.1).
For a hyperbolic relative cusp $q$, the fibre inertia already
injects in $P_U$. The tube map is consequently injective on
that fibre, and its image is $N_{E_U}(I_q)$: it induces the
identity on the base quotient, and both geometric kernels
are $I_q$ by peripheral self-normalization. The resulting
map of short exact sequences is an isomorphism. In frames
$e_j=g_{ij}e_i$ of $N$, root coordinates satisfy
$r_i=g_{ij}^{1/n}r_j$. Their triple-overlap defect is
$g_{ij}^{1/n}g_{jk}^{1/n}g_{ki}^{1/n}$, the positive Kummer
cocycle of $N$. This is exactly FLAGGED C6 and the sign in
Proposition~\ref{F-torsors}; dual conormal frames reverse it.

For the covering-germ criterion, more generally let
$H\to G$ be a homomorphism of profinite groups. If it is
injective and $H/K$ a finite transitive $H$-set, compactness
gives an open normal $M\lhd G$ with $M\cap H\subset K$.
Then $J=KM$ is open in $G$, $J\cap H=K$, and
$hK\mapsto hJ$ embeds $H/K$ into $G/J$ as an $H$-orbit.
Disjoint unions treat every finite $H$-set. Conversely a
nonidentity kernel element is detected by a finite quotient
of $H$ and acts trivially on every restricted $G$-set, so
the regular set of that quotient cannot embed. Taking
$H=T^J$ proves the asserted equivalence, including the
qualification ``sub-germ''.

Finite covers of a finite root stack are the covers of $U$
whose inertia orders divide the specified root orders. In
the chart $[\operatorname{Spec}A[w_i]/(w_i^{n_i}-z_i)/
\prod\mu_{n_i}]$, this is exactly unramifiedness on the
covering chart. Descent proves (F.43.2). Every finite
quotient of $P_U$ has finite inertia orders, so some
$M_{\bf n}$ lies in its kernel; hence their intersection
is trivial. For $(\mathbb A^1,0)$ the map before taking
roots is $T\xrightarrow{1}T$, whereas the finite-root
map is $T\to\mu_n$ with kernel $nT$. For a finite Kummer
\emph{cover}, the local lattice maps instead by a finite-index
inclusion $T'\hookrightarrow T^J$. A nonzero kernel $A$
persists because $mA\subset A\cap T'$ is nonzero whenever
$mT^J\subset T'$. These are the precise maps relating
the finite-root and covering operations.

For a punctured hyperbolic fibration, the geometric fibre
group is a centre-free free profinite group. The homotopy
boundary from the base's $\pi_2$ is central in that fibre
group and therefore vanishes. This gives the injected
fibre sequence of FLAGGED C4; a peripheral word injects
there by the full peripheral theorem. One can also verify
the completion step directly. For a discrete free fibre
$F$, conjugation maps the total group to
$\operatorname{Aut}(\widehat F)$; completion extends that
map, and its restriction to $\widehat F$ is the injective
inner-automorphism map. Thus the completed fibre injects.
A nontrivial cyclic subgroup of $F$ has its full profinite
topology: complete its immersed labelled circle to a finite
rose cover in which that circle is embedded, then use its
primitive integral homology class and pass to cores.

At a stable node $xy=t$, the contact map is
$(a,b)\mapsto a+b$, with kernel generated by $(1,-1)$;
this generator maps to the vanishing cycle in the fibre.
That cycle is essential: an inessential side would be an
unmarked rational tree attached at one point, violating
stability by summing the special-point inequalities over
its components. At a marking the kernel is its peripheral
loop; away from either locus it is zero. Compatible
toroidal charts give the corresponding finite-index or
primitive sublattices. The diagram of local lattices over
the injected fibre sequence therefore proves full inertia
injection by induction, starting with the stipulated
injective base lattices. Smith normal form gives
$\ker\widehat M=\widehat{\ker M}$ for the integer contact
matrix, so this argument proves the profinite assertion
as well as the discrete one. The transformations of
$z_i^{1/n}$ give precisely the claimed cyclotomic and
permutation actions and the contact matrix. The units
change tangential origins by Kummer cocycles.

For (F.43.3), substitution proves the equation and the
action is free for $s\ne0$. The normal loop maps to six
times the generator under the projection to the $t$-line,
so its full profinite map is injective. At infinity the
other inertia generator is the cusp loop of the punctured
elliptic fibre; the preceding injected fibre sequence
proves rank-two injectivity. For comparison the three
ordinary blowup charts of $f=y^2-x^3$ give
\[
 f=x^2(v^2-x),\qquad f=u^2v^3(v-u),\qquad
 f=u^6w^3(w-1).
\]
The exceptional tails have squares $-3,-2$ and their
normal loops $b,a$ satisfy $h=b^3=a^2$, where $h$ is
the third exceptional loop. The crossing maps are
$(m,n)\mapsto b^{m+3n}$ and $(m,n)\mapsto a^{m+2n}$.
Both $a$ and $b$ have full procyclic order, already
detected by their respective degrees $3$ and $2$ on the
$t$-line. The displayed kernels follow. Blowing up a
smooth boundary point of the product compactification of
the tripod squared likewise gives $(a,b)\mapsto a+b$
at its new crossing, with kernel $T(1,-1)$. This explains
exactly why arbitrary compactifications are not among
the recursively compatible stable models used above.

Finally, the full FLAGGED comparison can be made directly
on sections. C1 extends a rational geometric section
across a subset of codimension at least two: resolve its
graph; projective-space fibres of the resolution map
constantly to a positive-genus curve, so the graph
descends through the normal base. In genus zero an
\'etale-local boundary section makes its complement
affine, and the normal Hartogs equality extends the
coordinate functions. The inverse image of a remaining
boundary component would be a Cartier divisor supported
in codimension at least two, hence is empty. Repeat by
actual fibre products for a tower.

For C2 and C3, properness first extends a given trait
point to the compactified curve. If its boundary equation
pulls back as $z=t^\mu u$, the map on Kummer roots is
$z^{1/n}\mapsto t^{\mu/n}u^{1/n}$; its inertia map is
$T\xrightarrow{\mu}T\hookrightarrow P_U$. It is
nontrivial exactly when $\mu>0$. C5 combines this
codimension-one test with purity: the kernel of
$G_{K(S)}\to E_S$ is the closed normal subgroup generated
by divisorial inertias. C1 then extends across the
remaining codimension-two set. C8 is the injection of
normal frames into the splitting torsor proved in
Proposition~\ref{F-torsors}. C9 reduces the same root map
modulo $t$; it gives the normal-line map $v\mapsto u v^\mu$
and the specialization $\kappa(\bar u)$ of
Proposition~\ref{F-specialization}. Replacing $t$ by $at$
changes it by $-\mu\kappa(\bar a)$. Thus all the
comparison squares commute on each finite root cover
and therefore on their inverse limits. Each extension
test starts with an actual rational section and applies
these explicit maps; it supplies no missing point in
the open assertions of Propositions~\ref{F-field-reduction}
and~\ref{F-inductive-cuspidality}.
\end{proof}

\begin{proposition}[F.44 --- Proved: lifting at a potentially good base cusp]
\label{F-potential-boundary-lifting}
Let $K/\mathbb Q$ be finitely generated and let $U\to B$ be a
hyperbolic curve fibration over a hyperbolic curve. Let
$s:G_K\to E_U$ project to a cuspidal section $\tau$ at
$c\in(\overline B\setminus B)(K)$. Suppose the fibre over
$K((t))$, for a parameter $t$ at $c$, has potentially good
marked reduction. Let $Y_\tau/K$ be the smooth residue
curve obtained by the descent construction of
Proposition~\ref{F-reduction}; it has an identified arithmetic
extension
\[
 E_{Y_\tau}\simeq
 E_{U_{K((t))}}\mathbin{\times}_{G_{K((t))},\tau}G_K.
 \tag{F.44.1}
\]
Every $s$ whose induced section in (F.44.1) is a point or
cusp section has a birational lift $G_K\to G_{K(U)}$.
The lift can be obtained through the decomposition group
of a rank-two valuation of $K((t))(U)$ with residue $K$.

Consequently universal curve lifting implies surface
lifting for hyperbolic fibrations with potentially good
marked reduction at every rational base cusp after finite
constant extension. This extends the good-reduction range
of Proposition~\ref{F-good-boundary-lifting}; the implication
is an implication between the two universal assertions,
not an assertion of universal curve lifting.
\end{proposition}
\begin{proof}
Put $F=K((t))$ and $L_F=F(U_F)$. The comparison of relative
fibre sequences used in Proposition~\ref{F-good-boundary-lifting}
identifies $E_{U_F}$ with $E_U\times_{E_B}D_c$ and
$D_c$ with $G_F$. In particular the section $s$ has the
specified lift to the pullback along $\tau$. The exact
descent in Proposition~\ref{F-reduction} gives (F.44.1).
We first explain the valuation supplied by an actual
point or cusp of its residue curve, including compatibility
on finite coverings.

Choose a finite Galois extension $F_1/F$ over which the
stable marked model is smooth, and write $k_1$ for its
residue field. Its generic Galois isomorphisms extend
uniquely to that model. With $G=\operatorname{Gal}(F_1/F)$,
the homomorphism $\alpha:G_K\xrightarrow{\tau}G_F\to G$
defines exactly the descent datum used for $Y_\tau$.
After a finite residue extension killing $\alpha$, a
point or cusp of $Y_\tau$ is therefore a closed point
$y$ of the smooth special marked model, fixed by this
specified twisted descent datum.

At this point take the flag valuation: first the vertical
divisor, then the point on its smooth special curve.
It has value group $\mathbb Z\oplus_{\rm lex}\mathbb Z$
and a finite residue extension of $K$. For a cusp the
second parameter is its horizontal boundary equation;
for an interior point it is any parameter on the special
curve. Multiplying either parameter by a unit leaves
the valuation ring unchanged. Restriction through the
finite extension of function fields gives a rank-two
valuation $v$ of $L_F$, with finite residue over $K$.
Finite extension changes its value lattice only by finite
index.

The needed prolongation of $v$ to the maximal extension
unramified over $U_F$ is fixed by $s(G_K)$. This fact
requires the full comparison on covers, not just its
outer representation on the fibre group, so we specify
the construction. A finite \'etale cover of $U_F$,
after a further finite extension of $F_1$, has a smooth
marked model induced by the good-reduction comparison.
Indeed over the strict henselian good model its geometric
cover is the lift of the corresponding special-fibre
cover; the Kummer covers at its horizontal marks give
the compactification. In residue characteristic zero
all these ramification indices are invertible. Finite
presentation descends the model to a finite extension.
Refining two such choices by their common finite
extension and a common dominating cover makes the
specialization maps commute.

The section on $Y_\tau$ provides a compatible pro-point
on the completed special-fibre covering tower, by
Proposition~\ref{F-fixed-dictionary}. For an interior
point this is its ordinary fibre functor; for a cusp it
is the point fixed by its decomposition group. The
functors in the preceding paragraph carry its coordinates
to the points of the good special models of every
generic cover. The flag valuations at these points
restrict compatibly: the first valuation is the base
trait valuation, and the second is the discrete
point valuation on its smooth normalization. Their
compatibility can also be read on the local covers
$t^{1/a},z^{1/b}$; a change of unit preserves the
same valuation. Thus they define a prolongation
$\widetilde v$ of $v$ on the maximal unramified extension.
The equivariant comparison in (F.44.1) is constructed
from these same stable-model isomorphisms. Since the
chosen pro-point is fixed by the induced section,
$s(G_K)$ fixes $\widetilde v$.

Let $D_v\subset G_{L_F}$ be the decomposition group of
a prolongation to a separable closure. Restriction maps
$D_v$ onto the stabilizer of $\widetilde v$ in
$E_{U_F}$: at every finite Galois level a decomposition
group surjects onto the decomposition group below it,
and compactness gives the inverse-limit surjection.
Write $\overline D_v$ for its image. We have proved
$s(G_K)\subset\overline D_v$. If $k(v)/K$ is the
finite residue extension, the image of $D_v$ in $G_K$
is the subgroup $G_{k(v)}$, using the residue embedding
of algebraic constants. The inclusion of $s(G_K)$
therefore forces $k(v)=K$.

It remains to lift the specified section, rather than
merely its image subgroup. The full-rank valuation
sequence of Proposition~\ref{F-valuation-structure} is
\[
 1\longrightarrow I_v\simeq T^2\longrightarrow D_v
        \longrightarrow G_K\longrightarrow1,
 \tag{F.44.2}
\]
and it splits. Its image has kernel
$\overline I_v$, the image of $T^2$, over $G_K$.
The equal-characteristic-zero henselian calculation
applies to $L_F$ although it is not finitely generated
over $K$: it uses the coefficient-field embedding of
$K$, the finite-index rank-two value lattice, and
compatible roots of a basis of that lattice, or
equivalently the iterated henselian flag model.
These are the hypotheses used for the inertia and
splitting step of Proposition~\ref{F-valuation-structure}.
The geometric group $P_U$ is torsion-free: its exact
hyperbolic fibre sequence has torsion-free curve kernel
and quotient. Consequently $\overline I_v\subset P_U$
is torsion-free. For each prime $\ell$ it is a
torsion-free finitely generated quotient of
$\mathbb Z_\ell(1)^2$, hence free and a direct summand
quotient. Choosing these module splittings and taking
their product gives a continuous
$\widehat{\mathbb Z}$-linear section
\[
                       j:\overline I_v\longrightarrow I_v.
\]
Both modules have scalar cyclotomic Galois action, so
$j$ is $G_K$-equivariant.

Fix any splitting $r$ of (F.44.2). Relative to its
image, the prescribed section $s$ has the unique form
$s(g)=c(g)\overline r(g)$ for a continuous cocycle
$c:G_K\to\overline I_v$. Then
$g\mapsto j(c(g))r(g)$ is a continuous splitting of
(F.44.2) mapping to $s$. Composing it with
$G_{L_F}\to G_{K(U)}$, the restriction homomorphism
associated to the actual field embedding
$K(U)\hookrightarrow F(U_F)$, gives the required
birational lift. Compatibility of restriction with
the finite-cover functors proves that its image in
$E_U$ is the original section.

Finally universal curve lifting is equivalent to the
universal point-or-cusp assertion by
Proposition~\ref{F-cuspidalization}. Its application to
$B$ gives a point or cusp projection. A point projection
uses its actual fibre as in
Proposition~\ref{F-proper-polycurve-lifting}. A cusp
projection uses $Y_\tau$ and the construction just
proved. The remaining case is genuinely bad marked
reduction, for which Proposition~\ref{F-reduction}
excludes any such smooth residue curve even after
finite extension. No compatibility statement for that
case has been inferred here.
\end{proof}

\subsection{A global tangential relation using two open homomorphisms}

\begin{theorem}[F.45 --- Proved: intrinsic recognition by two open Hom maps]
\label{F-two-Hom-tangents}
Let $K$ be a number field, $U=\mathbb P^1_K\setminus\{0,1,\infty\}$,
$E=\pi_1(U)$, and $\Delta=\ker(E\to G_K)$. Fix an actual lifted
cusp inertia $I_0\subset\Delta$ above $0$ and a continuous section
$s:G_K\to N_E(I_0)$. Its class is induced by a nonzero
$K$-rational tangent vector at $0$ precisely when the following
entirely group-theoretic relation holds.

There exist an open subgroup $H\subset E$ with image $G_K$, a cusp
inertia $J$ of $H$, a continuous section
$\sigma:G_K\to N_H(J)$ of the arithmetic projection, and two
continuous open homomorphisms over $G_K$,
\[
                       p,q:H\longrightarrow E,
\]
such that
\[
 \begin{gathered}
 p|_J:J\xrightarrow{\sim} I_0,\qquad
 q(j)=p(j)^2\quad(j\in J),\\
             [p\circ\sigma]=[q\circ\sigma]=[s]
       \quad\text{modulo conjugation by }\Delta.
 \tag{F.45.1}
 \end{gathered}
\]
Here a cusp inertia of $H$ is an intersection
$H\cap aI_i a^{-1}$ for one of the three ambient peripheral
classes and an actual geometric conjugator $a$. The peripheral
classes themselves are characterized by the cyclotomic-normalizer
criterion. Thus all the quantifiers in (F.45.1) refer to the one
ambient arithmetic group, its open subgroups, their continuous
homomorphisms, and their sections; the test does not first declare
any source section geometric.

The same relation after restriction to some finite extension
$L/K$ is equivalent to (F.45.1) over $K$. For every specified
$t\in K^\times$, one can moreover choose $H,J,p,q$ so that the
common-image relation in (F.45.1), with $\sigma$ in that cusp
packet, has the unique image $[s_t]$, the tangent with coordinate
$t$ in the standard coordinate $z$ at $0$.

This supplies a uniform global relation recognizing the geometric
splittings in F(g)(5), using two open Hom maps. It does not replace
those maps by a word involving only the automorphisms of the
tripod or only the two entries of a Grothendieck--Teichm\"uller
pair. That more restrictive formula remains the separate question
in Proposition~\ref{F-open-tangents}.
\end{theorem}
\begin{proof}
We first prove necessity of tangentiality from the displayed
relation, without assuming a section conjecture. The open subgroup
$H$ is the arithmetic fundamental group of an actual connected
finite \emph{\'etale} cover $V\to U$. Its image onto $G_K$ makes
$V$ geometrically connected. The subgroup $J$ is the inertia of
an actual lifted cusp $c$ of $V$. Since $\sigma(G_K)$ normalizes
$J$, the underlying cusp is $K$-rational. These identifications
come from the finite-cover correspondence and the inertia
normalizer identity, so they do not reconstruct a curve from an
unclassified abstract group.

The relative open Hom theorem
(Proposition~\ref{F-open-Hom}, or
Proposition~\ref{A:relative-hom}) realizes $p$ and $q$, up to
their specified geometric inner ambiguity, by dominant
$K$-morphisms
\[
                          P,Q:V\longrightarrow U.
\]
They extend to morphisms of smooth proper completions. Because
$p(J)=I_0$ and $q(J)=I_0^2$ are nontrivial, both send $c$ to
$0$: a morphism extending to an interior point kills the local
inertia. The index of the image of cusp inertia is the positive
ramification degree. Thus $P$ has degree one locally at $c$, and
$Q$ has degree two locally there. In a $K$-rational uniformizer
$u$ at $c$, their completed local expressions are
\[
 P=a u(1+uA(u)),\qquad
 Q=b u^2(1+uB(u)),\qquad a,b\in K^\times.
 \tag{F.45.2}
\]
This also verifies the direction of the exponent-two comparison:
under the isomorphism $p|_J$, the actual map $q|_J$ is
multiplication by two, exactly as required in (F.45.1).

Let $\beta\in H^1(G_K,T)=\widehat{K^\times}$ be the parameter
of $\sigma$ relative to the tangent $u=1$. The compatible roots
of the units $1+uA(u)$ and $1+uB(u)$ have residue one. The local
root-cover calculation of Part B and
Proposition~\ref{F-specialization} therefore gives, additively,
\[
 [p\sigma]=\kappa_K(a)+\beta,
       \qquad [q\sigma]=\kappa_K(b)+2\beta
       \quad\text{in the packet at }0.
 \tag{F.45.3}
\]
Both images lie in that packet: normalizing $J$ normalizes its
nontrivial image, whose full arithmetic normalizer is
$N_E(I_0)$. Packet injectivity in $E$ converts the equality of
section classes in (F.45.1) into equality in (F.45.3). Subtraction
gives the exact formulas
\[
 \beta=\kappa_K(a/b),\qquad
                  [s]=\kappa_K(a^2/b).
 \tag{F.45.4}
\]
Thus the source parameter and the requested target parameter are
both ordinary $K$-rational tangential parameters. No division in
$\widehat{K^\times}$ and no local-to-global saturation claim is
needed; the two ramification degrees differ by one.

Conversely take $s=s_t$ for $t\in K^\times$. On the smooth open
curve
\[
 W=\mathbb P^1_K\setminus
       \bigl(\{0,1,\infty\}\cup\{z:z^2=t\}\bigr)
\]
there are the two dominant maps
\[
               P(z)=z,\qquad Q(z)=z^2/t
                       \quad\text{from }W\text{ to }U.
 \tag{F.45.5}
\]
Both have the cusp $0$, of local degrees one and two. The tangent
$u=z$ with value $t$ maps under the first to parameter $t$, and
under the second to $(1/t)t^2=t$. Its entire root-fibre-functor
section therefore has equal images $[s_t]$ under the two maps.

To place this source inside the one ambient group, use the marked
Belyi theorem over $K$. Precisely, Scherr--Zieve, Theorem 1.1
\cite{F-ScherrZieve}, gives a finite morphism
$b:\mathbb P^1_K\to\mathbb P^1_K$ with branch locus
$\{0,1,\infty\}$ and with every point of the finite complement
of $W$ mapped into that set. The theorem allows a finite
$G_K$-stable set of geometric points, so the roots of $z^2=t$
need not be $K$-rational. We opened the primary proof: its
construction is defined over $K$, and its additional ramification
condition at the prescribed points is stronger than what is
used here. Consequently
\[
 V=b^{-1}(U)\subset W,\qquad b:V\longrightarrow U
 \tag{F.45.6}
\]
is an actual geometrically connected finite \emph{\'etale} cover.
Deleting these additional points leaves the completed local ring
and the tangent at $0$ unchanged. Restrict (F.45.5) to $V$ and
use $b$ to identify $\pi_1(V)$ with an open subgroup $H\subset E$.
The two induced homomorphisms are open: each dominant map of
this particular pair has an explicit comparison. The first is
the open immersion $V\hookrightarrow U$. The second factors as
the open immersion
$V\hookrightarrow\mathbb P^1\setminus\{0,\infty,z^2=t\}$
followed by the finite \emph{\'etale} degree-two map $z\mapsto z^2/t$
to $U$. An open immersion of these normal curves induces a
surjection on fundamental groups, and the finite cover induces
an open injection. Both maps are over $G_K$. The cusp inertia $J$ and
the specified tangential section $\sigma$ now satisfy every part
of (F.45.1). Notice that the Belyi map may be ramified at $0$;
the local degrees one and two in the test concern $P,Q$, not
the inclusion of $H$ supplied by $b$.

For these fixed $H,J,p,q$, formula (F.45.3) has $a=1$ and
$b=1/t$. Equality forces $\beta=\kappa_K(t)$, hence has the
unique common image $[s_t]$. This proves the assertion for each
specified parameter and the global existential characterization
at once. Replacing any base paths conjugates the corresponding
subgroups and maps. If desired, the equality of section classes
can be made an equality of actual homomorphisms: a geometric
conjugator between two sections in this one packet normalizes
$I_0$ by cusp uniqueness, hence belongs to $I_0$ and acts
trivially on its inertia. Such adjustments preserve the two
inertia equations in (F.45.1).

We also verify descent for the germ version. Enlarge a finite
extension in which the test holds to a finite Galois extension
$L/K$. Relative to the standard tangent, let
$\alpha\in\widehat{K^\times}$ be the parameter of $s$.
The already proved direction over $L$ gives
$\operatorname{res}_{L/K}\alpha=\kappa_L(v)$ for $v\in L^\times$.
The left side is Galois-invariant. Injectivity of the ordinary
Kummer map over $L$ gives $\gamma(v)=v$ for every
$\gamma\in\operatorname{Gal}(L/K)$, so $v\in K^\times$.
Restriction
\[
 H^1(G_K,T)\longrightarrow H^1(G_L,T)
\]
is injective because $T^{G_L}=0$. Explicitly, a cocycle whose
restriction is zero after a coboundary adjustment has values
fixed by $G_L$, by the same cocycle identity used in
Proposition~\ref{F-centralizer}, and those values vanish.
Thus $\alpha=\kappa_K(v)$ and the test holds already over $K$
by the construction above.

For a small completely explicit Belyi witness, take $K=\mathbb Q$
and $t=2$. One can use
\[
 b(z)=2z^2-z^4,
 \qquad
 b^{-1}\{0,1,\infty\}
       =\{0,1,-1,\sqrt2,-\sqrt2,\infty\}.
\]
Its derivative is $4z(1-z^2)$, so its branch values are exactly
$0,1,\infty$. With $P=z$ and $Q=z^2/2$, this degree-four
cover supplies (F.45.1) and has unique common cusp image $[s_2]$.
The exact polynomial and local-degree checks are recorded in
\texttt{code/F/two\_hom\_tangents.py} and its output.

The manuscript anchor is cote 149, scan 85 and the specialization
construction of scans 85--95, together with F(g)(5) and the last
question of B(a). Our relation uses the full Hom theorem to
interpret open group maps and Belyi to put their sources inside
one arithmetic tripod. It invokes neither global section
surjectivity nor the stipulated local section theorem.
\end{proof}

\begin{proposition}[F.46 --- Proved: the finite-etale restriction selects exactly two rational tangents]
\label{F-finite-etale-tangents}
In Theorem~\ref{F-two-Hom-tangents} over $K=\mathbb Q$,
impose the additional requirement that both $p$ and $q$ be
injective. The resulting relation holds exactly for
\[
                         [s_{16}],\quad [s_{-16}].
 \tag{F.46.1}
\]
More generally, over a number field $K$, every tangent parameter
selected by this restricted relation belongs to $16\mu(K)$,
where $\mu(K)$ is the group of roots of unity in $K$.
This conclusion still holds when the local degrees one and two
in that relation are replaced by any two unequal positive
degrees $d,e$: require that $p$ and $q$ send the selected cusp
inertia into $I_0$ with these respective indices, and keep the
common-image equation for the section. Every common image is
then an ordinary tangent in $16\mu(K)$. In particular, over
$\mathbb Q$ even the union of all these unequal-degree
finite-\'etale tests selects exactly the two classes in
(F.46.1). The condition $d\ne e$ is essential: taking $p=q$
gives an equal-degree relation satisfied by every member of
the corresponding packet.

The additional requirement is exactly that the two reconstructed
morphisms from the common finite cover to the tripod are finite
\'etale. In that case $q\circ p^{-1}:p(H)\xrightarrow{\sim}q(H)$
is an actual isomorphism between open arithmetic subgroups.
Thus the forgetful map from these finite-cover relations to the
open-Hom relations of Theorem~\ref{F-two-Hom-tangents} has a
strictly smaller set of common section images: $[s_2]$ is an
explicit witness. This refutes replacing the two maps in that
criterion by such commensurations while retaining its inertia
degrees and equal-image equation. It makes no assertion that
every different formula using only GT automorphisms has this
same restriction.
\end{proposition}
\begin{proof}
We first justify the exact comparison between the group and
curve hypotheses. By the open Hom theorem, $p$ and $q$ are
induced by dominant maps $P,Q:V\to U$, where $V$ is the
finite \'etale cover attached to $H$ and
$U=\mathbb P^1\setminus\{0,1,\infty\}$. Suppose, for example,
that $p$ is injective. Its open image corresponds to a finite
\'etale cover $W\to U$, and $P$ lifts to $R:V\to W$ by the
finite-cover lifting criterion. The map on geometric fundamental
groups is an isomorphism. Both curves are affine and hyperbolic,
so their free profinite ranks give
$\chi(V)=\chi(W)<0$.

Extend $R$ to the smooth compactifications and let $e\ge1$
be its degree. Since the preimage of the boundary of $W$ is
contained in the boundary of $V$, Riemann--Hurwitz gives
\[
 \chi(V)=e\chi(W)-R_{\rm int}-m,
 \tag{F.46.2}
\]
where $R_{\rm int}\ge0$ is ramification over the interior of
$W$, and $m\ge0$ counts the additional deleted points of $V$
mapping to that interior. Equality of the Euler characteristics
makes $R_{\rm int}+m=(e-1)\chi(W)$, so $e=1$ and both
nonnegative terms vanish. Hence $R$ is an isomorphism and $P$
is finite \'etale. Conversely a finite \'etale map induces an
injective open fundamental-group map. The same argument applies
to $q$. This proves the asserted equivalence and the displayed
open-subgroup isomorphism; no general open homomorphism has
been treated as invertible.

We now compute the restriction on the local coefficients.
Fix an embedding $K\hookrightarrow\mathbb C$. Use the standard
uniformization
\[
 U(\mathbb C)=\overline\Gamma(2)\backslash\mathbb H,
 \qquad z=\lambda(\tau),\qquad
 \lambda(\tau)=16q_0+O(q_0^2),\quad q_0=e^{\pi i\tau}.
 \tag{F.46.3}
\]
Here $\overline\Gamma(2)$ is the image of the principal
level-two subgroup in $\operatorname{PSL}_2(\mathbb R)$.
The normalization, product expansion and identification of the
three punctures in (F.46.3) are recorded together in
\cite[\S1, equation (1.0.2), p.~628]{F-CalegariDimitrovTang};
the first coefficient also follows immediately from
$\lambda=\theta_2^4/\theta_3^4$ and
$\theta_2=2q_0^{1/4}+O(q_0^{9/4})$,
$\theta_3=1+2q_0+O(q_0^4)$.

Using $P$, identify the universal cover of $V(\mathbb C)$
with $\mathbb H$ and its deck group with a finite-index
subgroup $\Lambda\subset\overline\Gamma(2)$. The second
covering map has the form
\[
                  Q(\tau)=\lambda(\alpha\tau)
\]
for an orientation-preserving automorphism
$\alpha\in\operatorname{PSL}_2(\mathbb R)$ of $\mathbb H$.
Both $\Lambda$ and $\alpha\Lambda\alpha^{-1}$ have finite
index in $\overline\Gamma(2)$. Thus $\alpha$ commensurates
$\overline\Gamma(2)$.

For clarity, the required commensurator calculation is elementary.
The fixed points of its parabolic elements are exactly
$\mathbb P^1(\mathbb Q)$: an integral parabolic matrix has a
rational eigenspace, and each rational point is the image
of $\infty$ under an integral projective matrix, and a power
of the resulting parabolic lies in the finite-index level-two
group. Passing to any finite-index subgroup preserves this set,
since a positive power of every parabolic belongs to that
subgroup. Consequently $\alpha$ permutes
$\mathbb P^1(\mathbb Q)$. Its values at $0,1,\infty$ are
three distinct rational projective points; the unique projective
transformation having those three values has rational
coefficients. Hence
$\alpha\in\operatorname{PGL}_2^+(\mathbb Q)$.
Conversely any such rational matrix commensurates the integral
group: the lattices $\mathbb Z^2$ and $\alpha\mathbb Z^2$
are commensurable, and their simultaneous stabilizer has finite
index in each stabilizer. Intersecting with level two preserves
finite index. These are the actual two inclusions used in
the commensurator identification.

Let $c$ be the selected rational cusp of $V$. Since both
$P(c)$ and $Q(c)$ equal $0$, one may change the two universal
covering lifts by elements of $\overline\Gamma(2)$ so that
the chosen lift of $c$ is $\infty$ and
$\alpha(\infty)=\infty$. Therefore
\[
                \alpha\tau=a\tau+b,
                 \qquad a\in\mathbb Q_{>0},\quad b\in\mathbb Q.
\]
For the more general asserted local degrees $d,e$, the width
of the source cusp is $2d$, its analytic parameter is
$q_c=e^{\pi i\tau/d}$, and $a=e/d$. Writing
$\zeta=e^{\pi i b}$, a root of unity, gives
\[
                 P=16q_c^d+O(q_c^{2d}),\qquad
                 Q=16\zeta q_c^e+O(q_c^{2e}).
 \tag{F.46.4}
\]
If $u$ is instead a $K$-rational local parameter at $c$ and
$q_c=\gamma u+O(u^2)$, the leading coefficients are
$A=16\gamma^d$ and $B=16\zeta\gamma^e$. Their indicated
ratio is independent of that change of parameter:
\[
               A^e/B^d=16^{e-d}\zeta^{-d}.
 \tag{F.46.5}
\]
In particular $\xi=\zeta^{-d}$ belongs to $\mu(K)$, since
$A,B\in K^\times$. For the source packet parameter $\beta$
and the common target parameter $\alpha$, the exact local
Kummer maps are
\[
 \alpha=\kappa_K(A)+d\beta=\kappa_K(B)+e\beta,
 \qquad
 (e-d)\bigl(\alpha-\kappa_K(16)\bigr)=\kappa_K(\xi).
 \tag{F.46.6}
\]
This elimination uses actual homomorphisms of the packets; it
does not divide by $e-d$ in the power completion.

We check the torsion needed to extract the conclusion.
The divisor map gives an exact sequence from
$\mathcal O_K^\times/\mu(K)$ to $K^\times/\mu(K)$ and then
onto the group of principal fractional ideals. The first
group is free abelian of finite rank; the last is a subgroup
of the free abelian group on the finite primes, and is free
abelian as well. This exact sequence splits as a sequence of
abelian groups. Thus $K^\times/\mu(K)$ is free abelian, and
so is a chosen complement $F$ in
$K^\times\simeq\mu(K)\times F$. For a basis $(e_i)_{i\in I}$
of $F$, coordinate projections give an injective map
\[
       \widehat F=\varprojlim_n F/nF
              \longrightarrow\prod_{i\in I}\widehat{\mathbb Z}.
\]
Indeed a class with every coordinate zero has zero reduction
in the direct sum $F/nF$ for every $n$. The target is
torsion-free, and hence so is $\widehat F$. Completing the
displayed product decomposition therefore proves that the
torsion subgroup of $\widehat{K^\times}$ is exactly the image
of $\mu(K)$; this description is independent of the chosen
complement. Since $e-d\ne0$ and the right side of (F.46.6)
has finite order, $\alpha-\kappa_K(16)$ is torsion. It follows
that $\alpha\in\kappa_K(16\mu(K))$, as asserted. For the
original degrees $1,2$ the stronger literal coefficient
formula is $\alpha=\kappa_K(A^2/B)$.

For $K=\mathbb Q$ the only possibilities are $16$ and $-16$.
Changes of the universal paths conjugate the group maps; they
do not change either reconstructed rational function or the
local coefficient ratio. If $e=d$, the elimination instead
has zero left side and imposes no such restriction; the
identity example in the statement verifies this exception.

Both rational possibilities occur on the same explicit curve:
\[
 V=\mathbb P^1_{\mathbb Q}\setminus\{0,1,-1,\infty\},
 \qquad P(u)=\frac{4u}{(1+u)^2},\qquad
 Q_+(u)=u^2,\quad Q_-(u)=\frac{u^2}{u^2-1}.
 \tag{F.46.7}
\]
Each is a degree-two map with complete preimage of
$\{0,1,\infty\}$ equal to the displayed four-point boundary,
and all its ramification is there. Thus each restriction to
$V$ is finite \'etale. At $u=0$ their leading terms are
$4u,u^2,-u^2$, respectively. The source tangents $4$ and $-4$
have common images $s_{16}$ for $(P,Q_+)$ and $s_{-16}$ for
$(P,Q_-)$. Taking the image of the finite \'etale map $P$
as $H\subset E$ supplies the actual ambient-subgroup witnesses.
The rational functions, boundary preimages, ramification orders
and coefficients are checked in
\texttt{code/F/finite\_etale\_tangents.py} and its output.

Finally $s_2$ satisfies the unrestricted relation by the explicit
degree-four witness of Theorem~\ref{F-two-Hom-tangents}, while
$2\notin\{16,-16\}$. This proves strictness through the
specified forgetful map. The more general GT-only question
retained in Proposition~\ref{F-open-tangents} has not been
identified with this finite-etale restriction.
\end{proof}

\begin{proposition}[F.47 --- Proved: discontinuity on the cyclotomic kernel itself]
\label{F-cyclotomic-discontinuous-action}
Put $K=\mathbb Q(i)$, $\Gamma=G_K$,
$\Gamma^\circ=G_{K(\mu_\infty)}$, and
$U=\mathbb P^1_K\setminus\{0,1,\infty\}$.
There is an abstract semilinear $\Gamma$-action on the
geometric universal covering $\mathcal U$ of this same
tripod for which the following statements hold.
\begin{enumerate}
\item Its restriction to every open subgroup of
$\Gamma^\circ$ is not functionwise continuous. Its action
on the finite geometric homology quotient
$H_1(U_{\overline K},\mathbb F_3)$ is already discontinuous.
\item No open subgroup of $\Gamma^\circ$ admits a
finite-type quotient $\mathcal U\to W$ which is a
connected pro-\'etale covering and on which the induced
semilinear action is functionwise continuous. In particular,
no open subgroup of $\Gamma$ admits such a quotient.
\item The completed universal covering nevertheless has
exactly one fixed pro-point under this abstract
$\Gamma^\circ$-action, namely the chosen lift of $1/2$.
It is also fixed by the whole abstract $\Gamma$-action.
\end{enumerate}
Thus the inclusion of continuous actions into abstract
actions is strict already on the actual cyclotomic-kernel
domain occurring in (B$'$). The continuous finite-form
descent of Proposition~\ref{F-continuous-action} does not
extend to these actions. The universal fixed-point
existence assertion for arbitrary abstract actions is
neither proved nor refuted by this example.
\end{proposition}
\begin{proof}
Choose distinct rational primes $p_j\equiv1\pmod4$ and,
for each $j$, one Gaussian prime $\pi_j$ above $p_j$.
There are infinitely many such rational primes: from any
finite list take $N$ divisible by twice their product;
every odd prime divisor of $N^2+1$ is a new prime
congruent to one modulo four. The Euclidean domain
$\mathbb Z[i]$ gives
$p_j=\pi_j\overline\pi_j$ up to a unit, with the two
prime ideals distinct. Let
\[
 \epsilon_j:\Gamma\longrightarrow\mathbb F_2
\]
be the continuous quadratic character of
$K(\sqrt{\pi_j})/K$.

These characters remain linearly independent after
restriction to $\Gamma^\circ$. To prove this, a nonempty
finite product $a=\prod_j\pi_j^{e_j}$, with
$e_j\in\{0,1\}$, is not a square in $K$ by its prime
valuations. If its character vanished on
$\Gamma^\circ$, then
\[
 K(\sqrt a)\subset K(\mu_\infty)=\mathbb Q(\mu_\infty).
 \tag{F.47.1}
\]
Every finite subextension of the abelian extension
$\mathbb Q(\mu_\infty)/\mathbb Q$ is Galois over
$\mathbb Q$. Consequently complex conjugation would
preserve $K(\sqrt a)$ and its quadratic Kummer line
over $K$. This would give
$a/\overline a\in K^{\times2}$. At any selected
$\pi_j$ its valuation is instead one, a contradiction.
Thus every finite coordinate projection of
\[
 \epsilon=(\epsilon_j)_j:\Gamma^\circ
       \longrightarrow P=\prod_{j\ge1}\mathbb F_2
 \tag{F.47.2}
\]
is onto. Its compact image is closed, so (F.47.2) is
onto. The same coordinate map is also defined on all
of $\Gamma$ and is onto there.

Let $D=\bigoplus_{j\ge1}\mathbb F_2\subset P$.
Choose an abstract $\mathbb F_2$-linear functional
$\lambda:P/D\to\mathbb F_2$ taking the class of
$(1,1,\ldots)$ to one; extension of this nonzero vector
to a basis gives such a functional. With the quotient
by $D$ understood, define
\[
 \psi=\lambda\epsilon:\Gamma\longrightarrow\mathbb F_2.
 \tag{F.47.3}
\]
The restriction of $\psi$ to $\Gamma^\circ$ is nonzero
and its kernel is dense. Indeed it contains the inverse
image of the dense subgroup $D$, and a continuous
surjection of compact groups is an open quotient. The
kernel cannot be open, since a continuous functional on
$P$ vanishing on $D$ vanishes everywhere. The restriction
to every open subgroup $H\subset\Gamma^\circ$ is still
nonzero with dense kernel: its kernel is dense in $H$;
if it contained $H$, the full kernel would be an open,
hence closed, proper dense subgroup. Continuity on any
such $H$ would likewise make the full kernel open.

Choose a geometric pro-point $\widetilde x$ over
$x=1/2$, and let $s$ be the continuous arithmetic point
action of $\Gamma$ on $\mathcal U$. The involution
$a:z\mapsto1-z$ has a unique lift $\widetilde a$ fixing
$\widetilde x$. Existence is the lift of the based
finite-\'etale fibre functor; uniqueness follows because
a deck transformation fixing a point of a geometric
fibre is the identity. It follows that
$\widetilde a^2=1$ and that $\widetilde a$ commutes with
$s(\Gamma)$: conjugation by $s(g)$ gives another lift
of the same $K$-involution fixing $\widetilde x$.
Define the actual action
\[
 r(g)=\widetilde a^{\psi(g)}s(g),\qquad g\in\Gamma.
 \tag{F.47.4}
\]
The two commuting actions and the homomorphism (F.47.3)
prove the group law. Each $r(g)$ is an algebraic
semilinear automorphism, with the original action on
$\overline K$; normalization extends it to the completed
cover. On the coordinate $z$, its value is $z$ or
$1-z$ according to $\psi(g)$. Its stabilizer is therefore
$\ker\psi$, which is not open on any $H$ above.
For a finite quotient detecting the same failure, the
natural action on
$H_1(U_{\overline K},\mathbb F_3)$ is multiplication by
the cyclotomic character: the two independent classes
are the oriented inertia classes at $0$ and $1$.
It is trivial on $\Gamma^\circ$, whereas
$\widetilde a$ interchanges these two classes.
Consequently (F.47.4) acts on this quotient by the
nontrivial interchange to the power $\psi(g)$.

Here is the precise obstruction to a continuously
descended finite quotient. Let $\Omega$ be the covering
function field, let $H\subset\Gamma^\circ$ be open,
and suppose a finite-type quotient field
$M\subset\Omega$ is $r(H)$-stable with functionwise
continuous induced action. For every $f\in M$, the maps
\[
 H\longrightarrow\Omega,\qquad
 g\longmapsto r(g)f,\quad g\longmapsto s(g)f
\]
are continuous with $\Omega$ discrete. They agree on
the dense subgroup $H\cap\ker\psi$, hence everywhere.
Choose $g\in H$ with $\psi(g)=1$. Commutation in
(F.47.4) then gives $\widetilde a f=f$, so
$M\subset\Omega^{\widetilde a}$. If
$\mathcal U\to W$ were a connected pro-\'etale covering
with field $M$, the nontrivial $\widetilde a$ would be
a deck automorphism fixing $\widetilde x$. This is
impossible: unique lifting on finite connected
\'etale levels makes the deck action free on a geometric
fibre, and the inverse limit preserves that property.
This proves assertion 2 for $H$. If an open subgroup
of $\Gamma$ supplied such a quotient, its intersection
with $\Gamma^\circ$ would supply one for an open $H$,
giving the stated further conclusion.

Finally $\widetilde x$ is fixed by (F.47.4). Any pro-point
fixed by $r(\Gamma^\circ)$ is fixed by the natural
$s$-action of the dense subgroup
$\ker(\psi|_{\Gamma^\circ})$. Its stabilizer for the
natural action is closed: this holds at every finite
covering level, and then in their inverse limit.
It is therefore fixed by all $s(\Gamma^\circ)$.
Proposition~\ref{F-fixed-dictionary} says that this fixed
set is exactly $\{\widetilde x\}$. This proves assertion
3 through the actual comparison of the two actions.
Unlike the action in Proposition~\ref{F-discontinuous-action},
the twist here remains nontrivial and discontinuous on
$\Gamma^\circ$ itself, while the fixed-point assertion
continues to hold. No statement of universal existence
for other abstract actions follows from this comparison.
\end{proof}
